% !TEX root = ../main.tex

\chapter{مرور پژوهش‌ها و کارهای مرتبط}

% =========================================================
% 3-1
% =========================================================
\section{مقدمه}

توسعه یک سامانه گفت‌وگوی صوتی فارسی نیازمند بررسی چند حوزه
پژوهشی به‌صورت هم‌زمان است.
از یک سو، سامانه باید بتواند گفتار فارسی را با دقت و سرعت مناسب
به متن تبدیل کند؛ از سوی دیگر، مدل زبانی باید قادر به پردازش
زبان فارسی، درک درخواست کاربر و تولید پاسخ مناسب باشد.
در نهایت، سامانه تبدیل متن به گفتار باید پاسخ تولیدشده را
به‌صورت صوتی طبیعی و قابل فهم ارائه کند.

به همین دلیل، پژوهش‌های مرتبط با این پروژه را نمی‌توان تنها
در یک حوزه محدود بررسی کرد.
ادبیات مرتبط شامل پژوهش‌های انجام‌شده در زمینه تشخیص خودکار
گفتار فارسی، مدل‌های گفتاری چندزبانه، مدل‌های زبانی فارسی و
چندزبانه، تبدیل متن به گفتار فارسی و در نهایت سامانه‌های
مکالمه‌ای صوتی یکپارچه است.

هدف این فصل صرفاً معرفی معماری مدل‌ها نیست؛ زیرا مبانی نظری
آن‌ها در فصل قبل بررسی شده است.
در این فصل تلاش می‌شود پژوهش‌های پیشین از نظر مسئله مورد بررسی،
داده‌های مورد استفاده، روش پیشنهادی، معیارهای ارزیابی، نتایج
اصلی و محدودیت‌ها تحلیل شوند.

این بررسی همچنین به مشخص‌شدن جایگاه پروژه حاضر کمک می‌کند.
در پروژه حاضر، تمرکز تنها بر دستیابی به بهترین عملکرد مستقل
در یک مؤلفه نیست، بلکه هدف انتخاب و یکپارچه‌سازی مجموعه‌ای از
مدل‌های قابل اجرای محلی برای ساخت یک سامانه کامل

\begin{center}
\lr{ASR--\lr{LLM}--\lr{TTS}}
\end{center}

برای زبان فارسی است.

ازاین‌رو، یکی از موضوعات مهم در مرور پژوهش‌های قبلی،
توجه هم‌زمان به کیفیت، هزینه محاسباتی، سرعت استنتاج،
محدودیت منابع و امکان استفاده عملی از مدل‌ها در شرایط واقعی
است.

ساختار این فصل به این صورت است که ابتدا پژوهش‌ها و منابع
مربوط به تشخیص خودکار گفتار فارسی بررسی می‌شوند.
سپس مطالعات مرتبط با مدل‌های گفتاری چندزبانه و زبان‌های
کم‌منبع معرفی خواهند شد.
در ادامه، پژوهش‌های مربوط به مدل‌های زبانی فارسی و اجرای
محلی آن‌ها، مطالعات تبدیل متن به گفتار فارسی و سامانه‌های
گفت‌وگوی صوتی بررسی می‌شوند.
در پایان، یافته‌های پژوهش‌های پیشین با یکدیگر مقایسه شده و
شکاف پژوهشی که پروژه حاضر در پی پاسخ به آن است مشخص خواهد شد.


% =========================================================
% 3-2
% =========================================================
\section{پژوهش‌های مرتبط با تشخیص گفتار فارسی}

توسعه سامانه‌های تشخیص خودکار گفتار برای زبان فارسی با چالش‌هایی
مانند محدودیت مجموعه‌داده‌های بزرگ و متنوع، تفاوت لهجه‌ها،
شرایط ضبط متفاوت، گفتار محاوره‌ای و ویژگی‌های نوشتاری زبان
فارسی همراه است.

در سال‌های اخیر، در دسترس قرار گرفتن مجموعه‌داده‌های بزرگ‌تر
و مدل‌های ازپیش‌آموزش‌دیده چندزبانه باعث شده است رویکردهای
جدیدی برای توسعه سامانه‌های تشخیص گفتار فارسی مورد بررسی
قرار گیرند.

در ادامه، تعدادی از پژوهش‌ها و منابع مهم مرتبط با موضوع این
پروژه بررسی می‌شوند.


% ---------------------------------------------------------
\subsection{پایگاه داده DeepMine}

یکی از منابع مهم برای پژوهش‌های پردازش گفتار فارسی،
پایگاه داده
\lr{DeepMine}
است که توسط
\lr{Zeinali et al.}
معرفی شده است
\cite{zeinali2019deepmine}.

این مجموعه‌داده با هدف پشتیبانی از چند وظیفه مختلف پردازش
گفتار، از جمله تشخیص گوینده و تشخیص گفتار فارسی، توسعه یافته
است.

پایگاه داده DeepMine شامل داده‌های فارسی و انگلیسی از بیش از
۱۸۵۰ گوینده و در مجموع صدها هزار نمونه صوتی است.
بیش از ۴۸۰ ساعت از داده‌های صوتی این مجموعه دارای رونویسی
متنی هستند.

یکی از نقاط قوت این مجموعه، تنوع گویندگان از نظر سن،
جنسیت و لهجه است.
چنین تنوعی برای توسعه سامانه‌های تشخیص گفتاری که قرار است
در شرایط واقعی مورد استفاده قرار گیرند اهمیت زیادی دارد.

پژوهشگران علاوه بر معرفی داده‌ها، چند خط مبنا برای وظایف
تشخیص گوینده و تشخیص گفتار ارائه کردند.
نتایج آن‌ها نشان داد که این مجموعه می‌تواند برای آموزش و
ارزیابی مدل‌های مقاوم تشخیص گفتار فارسی مورد استفاده قرار گیرد.

\textbf{ارتباط با پروژه حاضر:}
وجود مجموعه‌هایی مانند DeepMine نشان می‌دهد که ارزیابی \lr{ASR}
فارسی باید تا حد امکان روی داده‌های متنوع از گویندگان و شرایط
ضبط مختلف انجام شود.
در پروژه حاضر نیز یکی از محدودیت‌های مهم، میزان تنوع مجموعه
ارزیابی است؛ موضوعی که در فصل بحث و تحلیل باید به آن توجه شود.


% ---------------------------------------------------------
\subsection{مجموعه‌داده \lr{Common Voice}}

یکی دیگر از منابع مهم مورد استفاده در پژوهش‌های تشخیص گفتار،
مجموعه
\lr{Mozilla Common Voice}
است \cite{ardila2020commonvoice}.

\lr{Common Voice} یک مجموعه‌داده چندزبانه گفتار است که داده‌های آن
به‌صورت جمع‌سپاری
\LTRfootnote{Crowdsourcing}
گردآوری و اعتبارسنجی می‌شوند.

هدف اصلی این پروژه فراهم‌کردن داده‌های عمومی گفتاری برای توسعه
فناوری‌های تشخیص گفتار در زبان‌های مختلف است.
در مقایسه با بسیاری از مجموعه‌داده‌های آزمایشگاهی، داده‌های
\lr{Common Voice} از تعداد زیادی گوینده جمع‌آوری می‌شوند و بنابراین
تنوع طبیعی بیشتری از نظر ویژگی‌های گوینده و شرایط ضبط دارند.

نسخه فارسی \lr{Common Voice} در چند پژوهش و همچنین در آموزش یا
ارزیابی مدل‌های فارسی استفاده شده است.

این مجموعه برای پروژه حاضر اهمیت ویژه‌ای دارد، زیرا بخش اصلی
ارزیابی مدل‌های تشخیص گفتار این پروژه بر داده‌های فارسی
\lr{Common Voice} انجام شده است.

\textbf{ارتباط با پروژه حاضر:}
استفاده از یک مجموعه آزمون مشترک برای مدل‌های مختلف امکان
مقایسه منصفانه‌تر آن‌ها را فراهم می‌کند.
با این حال، داده‌های \lr{Common Voice} الزاماً تمام شرایط یک
مکالمه طبیعی شامل نویز محیط، گفتار پیوسته، مکث‌ها و تعامل
چندمرحله‌ای را پوشش نمی‌دهند.
بنابراین نتایج حاصل از این مجموعه باید همراه با ارزیابی‌های
واقعی‌تر تفسیر شوند.


% ---------------------------------------------------------
\subsection{سازگارکردن wav2vec 2.0 برای گفتار فارسی کودکان}

یکی از پژوهش‌های مرتبط با استفاده از مدل‌های خودنظارتی برای
زبان فارسی توسط
\lr{Abaskohi et al.}
انجام شده است
\cite{abaskohi2022persianchildren}.

تمرکز این پژوهش بر تشخیص گفتار کودکان پیش‌دبستانی فارسی‌زبان
بود.
این مسئله دشوارتر از تشخیص گفتار بزرگسالان است، زیرا ویژگی‌های
آکوستیکی صدای کودکان، از جمله فرکانس و دامنه صوت، می‌تواند با
داده‌هایی که مدل‌های عمومی روی آن‌ها آموزش دیده‌اند تفاوت
قابل توجهی داشته باشد.

پژوهشگران از مدل
\lr{wav2vec 2.0}
به‌عنوان پایه استفاده کردند و یک هدف آموزشی جدید با عنوان

\begin{center}
\lr{Random Frequency Pitch (RFP)}
\end{center}

به فرآیند آموزش اضافه کردند.

هدف این روش افزایش مقاومت مدل نسبت به تغییر ویژگی‌های فرکانسی
گفتار بود.

مدل پیشنهادی علاوه بر مجموعه داده اختصاصی کودکان، بر بخش
فارسی \lr{Common Voice} نیز ارزیابی شد.
پژوهش گزارش کرد که روش پیشنهادی روی بخش فارسی \lr{Common Voice}
به نرخ خطای واژه حدود

\begin{center}
\lr{WER = 6.45}
\end{center}

دست یافته است.

این مطالعه نمونه‌ای از اهمیت
\textbf{سازگارکردن مدل با دامنه}
\LTRfootnote{Domain Adaptation}
است.
حتی اگر یک مدل ازپیش‌آموزش‌دیده عمومی عملکرد مناسبی داشته باشد،
ویژگی‌های خاص جمعیت یا محیط هدف می‌توانند نیازمند تنظیم مجدد
مدل باشند.

\textbf{ارتباط با پروژه حاضر:}
نتایج این پژوهش نشان می‌دهد که عملکرد یک مدل \lr{ASR} تنها ویژگی
ذاتی معماری آن نیست و به داده‌های تنظیم و دامنه کاربرد نیز
وابسته است.
ازاین‌رو، مقایسه مدل‌های پروژه حاضر باید روی یک مجموعه داده
یکسان و با شرایط پردازشی یکسان انجام شود.


% ---------------------------------------------------------
\subsection{معیار PSRB برای تشخیص گفتار فارسی}

یکی از پژوهش‌های جدیدتر و مرتبط با ارزیابی جامع سامانه‌های
تشخیص گفتار فارسی،
\lr{Persian Speech Recognition Benchmark (PSRB)}
است \cite{sedghiyeh2025psrb}.

این پژوهش با هدف ایجاد یک معیار ارزیابی متنوع‌تر برای سامانه‌های
\lr{ASR} فارسی انجام شده است.
در این مطالعه، ده سامانه تشخیص گفتار شامل مدل‌های متن‌باز و
سامانه‌های تجاری مورد ارزیابی قرار گرفته‌اند.

یکی از نکات مهم PSRB توجه به تنوع شرایط زبانی و آکوستیکی است.
پژوهشگران تنها روی گفتار استاندارد تمرکز نکرده‌اند و عملکرد
سامانه‌ها در شرایطی مانند لهجه‌های منطقه‌ای و گفتار کودکان
را نیز بررسی کرده‌اند.

نتایج گزارش‌شده نشان می‌دهد که حتی مدل‌هایی که روی فارسی
استاندارد عملکرد مناسبی دارند، در برخی شرایط مانند
لهجه‌های منطقه‌ای، گفتار کودکان و چالش‌های زبانی خاص با افت
عملکرد مواجه می‌شوند.

علاوه بر \lr{WER} استاندارد، این پژوهش تحلیل دقیق‌تری از انواع
خطاهای رونویسی ارائه کرده و معیاری پیشنهاد کرده است که به
خطاهای جایگزینی وزن متفاوتی می‌دهد تا خطاهای جزئی و
جایگزینی‌های نزدیک به پاسخ صحیح بهتر تحلیل شوند.

\textbf{ارتباط با پروژه حاضر:}
این مطالعه نشان می‌دهد که یک عدد \lr{WER} واحد برای توصیف کامل
عملکرد یک سامانه تشخیص گفتار کافی نیست.
در پروژه حاضر نیز علاوه بر \lr{WER} از \lr{CER}، زمان تأخیر،
\lr{RTF} و مصرف منابع برای مقایسه مدل‌ها استفاده می‌شود.
در مراحل تکمیلی، بررسی مدل منتخب در شرایط نویزی و گفتار
واقعی‌تر می‌تواند ارزیابی را به شرایط مطرح‌شده در PSRB
نزدیک‌تر کند.


% ---------------------------------------------------------
\subsection{پردازش زبان‌های دارای خط فارسی و عربی}

پژوهش
\lr{Bandarupalli et al.}
بر چالش‌های پردازش زبان‌هایی با خط فارسی--عربی در سامانه‌های
تشخیص گفتار تمرکز کرده است
\cite{bandarupalli2025persoarabic}.

این پژوهش زبان‌های فارسی، عربی و اردو را بررسی کرده و نشان
می‌دهد که مراحل پیش‌پردازش متن، توکن‌سازی و تبدیل واجی نقش
مهمی در طراحی سامانه \lr{ASR} برای این زبان‌ها دارند.

پژوهشگران یک خط لوله پردازشی شامل پاک‌سازی داده،
توکن‌سازی و واج‌سازی
\LTRfootnote{\lr{Phonemization}}
ارائه کردند و مراحل آن را با کمک ارزیابی متخصصان زبان اصلاح
کردند.

این موضوع برای زبان فارسی اهمیت دارد، زیرا تفاوت نویسه‌های
فارسی و عربی، نیم‌فاصله، اشکال مختلف نوشتاری و نحوه نرمال‌سازی
متن می‌تواند مستقیماً بر آموزش و حتی محاسبه \lr{WER} اثر بگذارد.

\textbf{ارتباط با پروژه حاضر:}
پیش از محاسبه \lr{WER} و \lr{CER} باید متن مرجع و خروجی مدل با قواعد
یکسانی نرمال‌سازی شوند.
در غیر این صورت، بخشی از اختلاف گزارش‌شده ممکن است ناشی از
تفاوت نوشتاری باشد نه خطای واقعی در تشخیص محتوای گفتار.


% ---------------------------------------------------------
\subsection{مدل‌های کوچک‌تر برای زبان‌های کم‌منبع}

در مطالعه‌ای دیگر،
\lr{Bandarupalli et al.}
استفاده از پیش‌آموزش پیوسته چندزبانه برای زبان‌های کم‌منبع
را بررسی کردند
\cite{bandarupalli2025efficient}.

در این پژوهش، فارسی همراه با عربی و اردو به‌عنوان نمونه‌هایی
از زبان‌های دارای منابع محدودتر مورد بررسی قرار گرفت.

پژوهشگران یک مجموعه چندزبانه شامل حدود ۳۰۰۰ ساعت گفتار
بدون برچسب ایجاد کردند و از پیش‌آموزش پیوسته همراه با
توکن‌سازی متناسب با ساختار زبانی استفاده کردند.

مدل پیشنهادی آن‌ها حدود ۳۰۰ میلیون پارامتر داشت و طبق نتایج
گزارش‌شده توانست روی فارسی عملکرد رقابتی ارائه کند و در
ارزیابی انجام‌شده از
\lr{Whisper Large-v3}
بهتر عمل کند، با وجود آنکه اندازه مدل به‌مراتب کوچک‌تر بود.

\textbf{ارتباط با پروژه حاضر:}
این نتیجه از یکی از فرض‌های طراحی پروژه حاضر پشتیبانی می‌کند:
بزرگ‌ترین مدل الزاماً بهترین مدل برای یک سامانه محلی نیست.
کیفیت داده، سازگاری مدل با زبان هدف و هزینه محاسباتی باید
به‌صورت هم‌زمان مورد توجه قرار گیرند.


% ---------------------------------------------------------
\subsection{گفتار خودانگیخته فارسی و سازگارکردن \lr{Whisper}}

یکی از چالش‌های مهم در تشخیص گفتار فارسی، تفاوت میان گفتار
خوانده‌شده و گفتار طبیعی یا خودانگیخته است.

\lr{Namdarzadeh and Ballier}
در سال ۲۰۲۶ مجموعه داده
\textbf{Spontaneous Persian \lr{Speech} (SPS)}
را برای مطالعه این مسئله معرفی کردند
\cite{namdarzadeh2026sps}.

این مجموعه حدود ۶۹۴ دقیقه گفتار از ۶۵ گوینده را شامل می‌شود
و داده‌ها در محیط‌های مختلف مانند خانه، خودرو، فروشگاه و
محیط‌های آرام و شلوغ ضبط شده‌اند.

پژوهشگران از این مجموعه برای تنظیم نسخه‌های
\lr{small}
و
\lr{medium}
مدل \lr{Whisper} استفاده کردند و کاهش قابل توجه \lr{WER} پس از
تنظیم روی گفتار خودانگیخته فارسی گزارش کردند.

اهمیت این پژوهش در این است که بسیاری از مجموعه‌های استاندارد
\lr{ASR} شامل جملات خوانده‌شده یا قطعات صوتی نسبتاً کنترل‌شده هستند،
در حالی که ورودی یک چت‌بات صوتی معمولاً ماهیت مکالمه‌ای و
خودانگیخته دارد.

\textbf{ارتباط با پروژه حاضر:}
از آنجا که هدف پروژه حاضر تعامل صوتی طبیعی با کاربر است،
ارزیابی صرف روی داده‌های خوانده‌شده نمی‌تواند تمام رفتار
سامانه را نشان دهد.
در نتیجه، در کنار بنچمارک اصلی، ایجاد یک مجموعه آزمایش کوچک
از گفتار طبیعی و مکالمه‌ای می‌تواند ارزیابی نهایی پروژه را
واقع‌بینانه‌تر کند.


% ---------------------------------------------------------
\subsection{تأثیر قطعه‌بندی صوت بر عملکرد \lr{Whisper}}

پژوهش دیگری در سال ۲۰۲۶ به‌طور خاص تأثیر روش
قطعه‌بندی صوت را بر عملکرد \lr{Whisper} در زبان فارسی بررسی کرده
است
\cite{iranmanesh2026segmentation}.

در این مطالعه، حدود ۱۰ ساعت محتوای فارسی از YouTube مورد
استفاده قرار گرفته و دو روش قطعه‌بندی مقایسه شده‌اند:
قطعه‌بندی بر اساس زمان‌بندی متن مرجع و قطعه‌بندی مبتنی بر
سکوت.

نتایج نشان دادند که روش قطعه‌بندی مناسب وابسته به نوع محتوا
است.
در محتوای پادکست، قطعه‌بندی مبتنی بر زمان‌بندی مرجع \lr{WER}
پایین‌تری ایجاد کرده، در حالی که در محتوای سرگرمی
قطعه‌بندی مبتنی بر سکوت عملکرد بهتری داشته است.

این نتیجه نشان می‌دهد که نحوه تعیین مرزهای گفتار می‌تواند
به اندازه انتخاب خود مدل در عملکرد \lr{ASR} مؤثر باشد.

\textbf{ارتباط با پروژه حاضر:}
این یافته مستقیماً به طراحی VAD و تشخیص پایان گفتار در
چت‌بات حاضر مرتبط است.
اگر جمله کاربر خیلی زود یا خیلی دیر قطعه‌بندی شود،
ممکن است دقت \lr{ASR} کاهش یابد.
بنابراین VAD در پروژه حاضر نباید صرفاً به‌عنوان یک ابزار
برای کاهش سکوت در نظر گرفته شود، بلکه بخشی از فرآیند
بهینه‌سازی عملکرد \lr{ASR} است.


% ---------------------------------------------------------
\subsection{مدل \lr{FastConformer} فارسی NVIDIA}

در کنار مقالات پژوهشی، مدل‌های ازپیش‌آموزش‌دیده اختصاصی فارسی
نیز نقش مهمی در توسعه سامانه‌های کاربردی دارند.

مدل

\begin{center}
\lr{nvidia/stt\_fa\_fastconformer\_hybrid\_large}
\end{center}

یک مدل تشخیص گفتار فارسی مبتنی بر
\lr{FastConformer}
است که توسط NVIDIA ارائه شده است
\cite{nvidiaPersianFastConformer}.

این مدل حدود ۱۱۵ میلیون پارامتر دارد و از ساختار ترکیبی
\lr{Transducer--CTC}
استفاده می‌کند.
مدل با استفاده از داده‌های فارسی
\lr{Mozilla Common Voice 15.0}
تنظیم شده است.

طبق مستندات مدل، رمزگذار از معماری \lr{FastConformer} استفاده
می‌کند و دو سر خروجی CTC و Transducer در دسترس هستند.

این مدل برای پروژه حاضر اهمیت ویژه‌ای دارد، زیرا یکی از
مدل‌های اصلی مورد ارزیابی در آزمایش‌ها است.
با این حال، نتایج گزارش‌شده توسط ارائه‌دهنده مستقیماً با
نتایج پروژه حاضر قابل مقایسه نیستند، زیرا نحوه تقسیم داده،
روش نرمال‌سازی و تنظیمات رمزگشایی ممکن است متفاوت باشند.

به همین دلیل، در این پروژه مدل \lr{FastConformer} همراه با سایر
مدل‌ها روی مجموعه آزمون و خط لوله ارزیابی یکسان مجدداً
بنچمارک شده است.


% ---------------------------------------------------------
\subsection{جمع‌بندی مطالعات \lr{ASR} فارسی}

مرور پژوهش‌های فوق چند نکته مهم را درباره توسعه تشخیص گفتار
فارسی نشان می‌دهد.

نخست، کیفیت و تنوع داده نقش اساسی در عملکرد مدل دارند.
مجموعه‌هایی مانند DeepMine، \lr{Common Voice}، PSRB و SPS تلاش
کرده‌اند بخش‌های متفاوتی از تنوع گفتار فارسی را پوشش دهند.

دوم، مدل‌های ازپیش‌آموزش‌دیده چندزبانه مانند wav2vec 2.0
و \lr{Whisper} امکان توسعه سامانه‌های فارسی را با نیاز کمتر به
آموزش از ابتدا فراهم کرده‌اند، اما عملکرد آن‌ها به میزان
سازگاری داده‌های آموزش با شرایط کاربرد وابسته است.

سوم، نتایج پژوهش‌های اخیر نشان می‌دهند عواملی مانند
نرمال‌سازی متن، نوع توکن‌سازی، قطعه‌بندی صوت، سن گوینده،
لهجه و شرایط آکوستیکی می‌توانند اثر قابل توجهی بر \lr{WER} داشته
باشند.

در نهایت، برای یک سامانه گفت‌وگوی صوتی محلی، دقت تنها معیار
انتخاب مدل نیست.
سرعت، میزان حافظه، قابلیت اجرای محلی و رفتار مدل در گفتار
واقعی نیز باید در تصمیم‌گیری دخیل باشند.

این موضوع مبنای روش پروژه حاضر برای مقایسه چند مدل \lr{ASR}
بر اساس مجموعه‌ای از معیارهای دقت و کارایی محاسباتی است.

% =========================================================
% 3-3
% =========================================================
\section{مدل‌های چندزبانه و زبان‌های کم‌منبع برای تشخیص گفتار}

یکی از چالش‌های اصلی توسعه سامانه‌های تشخیص گفتار برای زبان فارسی،
محدودیت نسبی داده‌های گفتاری برچسب‌خورده در مقایسه با زبان‌هایی
مانند انگلیسی است.
این مسئله سبب شده است روش‌های
\textbf{انتقال میان‌زبانی}
\LTRfootnote{Cross-Lingual Transfer}
و
\textbf{پیش‌آموزش چندزبانه}
\LTRfootnote{Multilingual Pre-training}
به یکی از رویکردهای مهم در توسعه سامانه‌های \lr{ASR} برای زبان‌های
کم‌منبع تبدیل شوند.

ایده اصلی این روش‌ها آن است که مدل پیش از تنظیم برای زبان هدف،
از داده‌های صوتی چندین زبان برای یادگیری بازنمایی‌های عمومی
گفتار استفاده کند.
در نتیجه، بخشی از ساختارهای آکوستیکی و واجی آموخته‌شده از
زبان‌های پرمنبع می‌تواند برای زبان‌هایی با داده محدود نیز مفید
باشد.

این رویکرد برای پروژه حاضر اهمیت ویژه‌ای دارد، زیرا چند مورد از
مدل‌های بررسی‌شده، از جمله
\lr{wav2vec2/XLSR}
و
\lr{Whisper}،
به‌صورت چندزبانه پیش‌آموزش دیده‌اند.


% ---------------------------------------------------------
\subsection{انتقال میان‌زبانی در XLSR}

یکی از پژوهش‌های بنیادی در این حوزه توسط
\lr{Conneau et al.}
ارائه شده است
\cite{conneau2020xlsr}.

در این پژوهش، مدل
\lr{XLSR}
با هدف یادگیری بازنمایی‌های مشترک گفتار از چند زبان مختلف
توسعه یافت.
مدل بر پایه wav2vec 2.0 ساخته شده و به‌جای پیش‌آموزش روی تنها
یک زبان، از داده‌های خام صوتی چندین زبان استفاده می‌کند.

پژوهشگران نشان دادند که پیش‌آموزش چندزبانه می‌تواند در بسیاری
از زبان‌ها نسبت به پیش‌آموزش تک‌زبانه عملکرد بهتری ایجاد کند،
به‌ویژه در شرایطی که داده برچسب‌خورده زبان هدف محدود است.

نسخه
\lr{XLSR-53}
با استفاده از ۵۳ زبان پیش‌آموزش داده شده است و هدف اصلی آن
ایجاد یک مدل مشترک بود که بتواند پس از تنظیم روی داده محدود
هر زبان، برای \lr{ASR} همان زبان استفاده شود.

نتایج این پژوهش نشان داد که بازنمایی‌های آموخته‌شده میان
زبان‌ها تا حدی مشترک هستند و میزان اشتراک در زبان‌های نزدیک‌تر
بیشتر است.
این موضوع نشان می‌دهد که ساختارهای آکوستیکی و واجی مشترک میان
زبان‌ها می‌توانند به انتقال دانش کمک کنند.

\textbf{ارتباط با پروژه حاضر:}
مدل
\lr{wav2vec2-large-xlsr-persian-v3}
که در این پروژه ارزیابی شده است، از همین ایده استفاده می‌کند.
بنابراین عملکرد آن را می‌توان نتیجه ترکیب پیش‌آموزش چندزبانه و
تنظیم اختصاصی برای زبان فارسی در نظر گرفت.

با این حال، نتایج پروژه حاضر نشان می‌دهند که استفاده از یک
مدل چندزبانه به‌تنهایی تضمین‌کننده بهترین عملکرد نیست و مدل‌های
اختصاصی‌تر فارسی نیز باید مورد مقایسه قرار گیرند.


% ---------------------------------------------------------
\subsection{محدودیت پیش‌آموزش چندزبانه}

با وجود مزایای مدل‌های چندزبانه، پژوهش‌های جدیدتر نشان داده‌اند
که عملکرد آن‌ها برای تمام زبان‌ها یکسان نیست.

\lr{San et al.}
در مطالعه‌ای درباره انتقال مثبت در \lr{ASR} کم‌منبع نشان دادند که
حتی مدل‌های بزرگ چندزبانه مانند
\lr{XLS-R}
ممکن است برای زبان‌هایی که در داده پیش‌آموزش حضور محدودی
داشته‌اند عملکرد ضعیف‌تری داشته باشند
\cite{san2024positive}.

در این پژوهش، پیش‌آموزش پیوسته
\LTRfootnote{Continued Pre-training}
روی مقدار محدودی از صوت بدون برچسب زبان هدف بررسی شد.
نتایج نشان دادند که افزودن داده از یک زبان دیگر که از نظر
زبانی یا آکوستیکی به زبان هدف نزدیک باشد، می‌تواند عملکرد را
بهبود دهد.

برای مثال، در برخی آزمایش‌ها تنها مقدار کمی داده از زبان
کم‌منبع همراه با داده یک زبان اهداکننده
\LTRfootnote{Donor Language}
تقریباً به اندازه استفاده از حجم بسیار بیشتری از داده زبان
هدف مؤثر بوده است.

\textbf{نتیجه مهم این پژوهش} آن است که در انتقال میان‌زبانی،
تنها حجم داده اهمیت ندارد و میزان شباهت میان زبان مبدأ و هدف
نیز می‌تواند نقش تعیین‌کننده داشته باشد.

\textbf{ارتباط با پروژه حاضر:}
این نتیجه نشان می‌دهد که عملکرد یک مدل چندزبانه فارسی باید
بر اساس داده واقعی فارسی ارزیابی شود و نمی‌توان تنها بر اساس
اندازه مدل یا تعداد زبان‌های پیش‌آموزش‌دیده درباره کیفیت آن
قضاوت کرد.


% ---------------------------------------------------------
\subsection{\lr{Whisper} در زبان‌های کم‌منبع}

مدل
\lr{Whisper}
یکی از مهم‌ترین مدل‌های چندزبانه تشخیص گفتار در سال‌های اخیر
است.
این مدل با استفاده از حجم بسیار بزرگی از داده‌های صوتی و متنی
چندزبانه آموزش داده شده است
\cite{radford2022whisper}.

یکی از نقاط قوت \lr{Whisper} توانایی استفاده مستقیم روی زبان‌های
متعدد بدون نیاز به آموزش مدل از ابتدا است.
با این حال، پژوهش‌های انجام‌شده در حوزه زبان‌های کم‌منبع نشان
می‌دهند که عملکرد اولیه مدل می‌تواند با تنظیم اختصاصی روی
داده زبان هدف بهبود قابل توجهی پیدا کند.

برای مثال،
\begin{latin}Pineiro-Martin et al.\end{latin}
در پژوهشی روی یادگیری چندزبانه پیوسته، \lr{Whisper} را برای مجموعه‌ای
از زبان‌های پرمنبع و یک زبان کم‌منبع تنظیم کردند
\cite{pineiro2024weighted}.

آن‌ها از یک تابع زیان وزن‌دهی‌شده برای کاهش اثر عدم‌توازن میان
زبان‌ها استفاده کردند.
نتایج نشان داد این روش در زبان کم‌منبع باعث کاهش
\lr{WER}
نسبت به مدل تنظیم‌شده معمولی و همچنین نسبت به \lr{Whisper} اصلی شد.

این یافته نشان می‌دهد که مدل‌های بسیار بزرگ چندزبانه نیز ممکن
است بدون سازگار‌سازی مناسب، برای زبان‌های کم‌منبع بهترین
عملکرد ممکن را ارائه ندهند.


% ---------------------------------------------------------
\subsection{روش‌های سازگاری \lr{Whisper} بدون تنظیم کامل}

در پژوهشی دیگر،
\lr{Hsu et al.}
روش
\lr{Meta-Whisper}
را برای بهبود \lr{ASR} زبان‌های کم‌منبع معرفی کردند
\cite{hsu2024metawhisper}.

هدف این روش کاهش نیاز به تنظیم کامل مدل بود.
پژوهشگران از نمونه‌های گفتاری مشابه و روش‌های انتخاب مبتنی بر
نزدیک‌ترین همسایه برای فراهم‌کردن نمونه‌های زمینه‌ای استفاده
کردند.

نتایج روی مجموعه
\lr{ML-SUPERB}
نشان داد که استفاده از نمونه‌های مناسب می‌تواند نرخ خطای نویسه
را برای زبان‌های کم‌منبع کاهش دهد.

این خط پژوهش نشان می‌دهد که علاوه بر تنظیم مستقیم وزن‌های مدل،
روش‌های دیگری نیز برای سازگارکردن مدل‌های بزرگ گفتاری با زبان
هدف وجود دارند.


% ---------------------------------------------------------
\subsection{اهمیت داده مرتبط با دامنه}

یکی از مهم‌ترین یافته‌های اخیر در \lr{ASR} زبان‌های کم‌منبع آن است
که
\textbf{مرتبط‌بودن داده با زبان و دامنه هدف می‌تواند از اندازه
مدل مهم‌تر باشد}.

در پژوهش
\lr{Bandarupalli et al.}
روی فارسی، عربی و اردو، یک مدل حدود ۳۰۰ میلیون پارامتری با
استفاده از پیش‌آموزش پیوسته روی یک مجموعه ۳۰۰۰ ساعته چندزبانه
توسعه داده شد
\cite{bandarupalli2025efficient}.

پژوهشگران سه راهبرد مختلف را بررسی کردند:
آموزش از ابتدا، ادامه پیش‌آموزش از XLS-R و ادامه پیش‌آموزش
از wav2vec 2.0 Large.

در نتایج گزارش‌شده برای فارسی، بهترین مدل آن‌ها به
\lr{WER}
حدود 17.1 درصد دست یافت، در حالی که \lr{Whisper Large-v3} در همان
ارزیابی حدود 21.4 درصد گزارش شد.
این در حالی بود که مدل پیشنهادی تنها حدود یک‌پنجم اندازه
\lr{Whisper Large-v3} را داشت.

این نتیجه از نظر طراحی سامانه‌های محلی بسیار مهم است.
مدل کوچک‌تر در صورتی که با داده مناسب‌تر آموزش دیده باشد،
می‌تواند نسبت به مدل بسیار بزرگ‌تر عملکرد بهتری ارائه کند.

\textbf{ارتباط با پروژه حاضر:}
یکی از یافته‌های اصلی بنچمارک پروژه حاضر نیز همین مسئله را
بررسی می‌کند.
انتخاب مدل نهایی نباید تنها بر اساس اندازه یا شهرت مدل انجام
شود، بلکه دقت روی فارسی، سرعت و مصرف منابع باید به‌طور مستقیم
اندازه‌گیری شوند.


% ---------------------------------------------------------
\subsection{مدل‌های چندزبانه در شرایط بسیار کم‌منبع}

در حوزه زبان‌هایی که حتی چند ده ساعت داده گفتاری برچسب‌خورده
در دسترس نیست، استفاده از مدل‌های چندزبانه اهمیت بیشتری پیدا
می‌کند.

پژوهش‌های جدید روی مدل‌هایی مانند
\lr{XLS-R}
و
\lr{MMS}
نشان داده‌اند که تنظیم مدل‌های ازپیش‌آموزش‌دیده می‌تواند در
زبان‌هایی با مقدار بسیار محدود داده نیز مفید باشد.

با این حال، میزان بهبود به عواملی مانند مقدار داده،
شباهت زبان هدف به زبان‌های موجود در پیش‌آموزش،
کیفیت رونویسی و نوع توکن‌سازی وابسته است.

این مطالعات نشان می‌دهند که در \lr{ASR} کم‌منبع نمی‌توان یک مدل
واحد را برای تمام زبان‌ها به‌عنوان بهترین انتخاب در نظر گرفت.


% ---------------------------------------------------------
\subsection{مدل‌های بسیار بزرگ در برابر مدل‌های کارآمد}

روند توسعه مدل‌های گفتاری در سال‌های اخیر از یک سو به سمت
مدل‌های بسیار بزرگ چندزبانه حرکت کرده است، اما از سوی دیگر
پژوهش‌های متعددی بر کاهش هزینه محاسباتی تمرکز کرده‌اند.

در یک سامانه ابری، استفاده از مدل چندمیلیارد پارامتری ممکن است
قابل قبول باشد، اما در یک سامانه محلی محدودیت حافظه و سرعت
بسیار مهم‌تر می‌شود.

بنابراین، برای یک کاربرد محلی می‌توان مسئله انتخاب مدل را به
صورت یک بهینه‌سازی چندهدفه در نظر گرفت:

\begin{equation}
M^*
=
\arg\min_M
\left(
\alpha E_M
+
\beta T_M
+
\gamma R_M
\right)
\end{equation}

که در آن:

\begin{itemize}
    \item $E_M$ خطای تشخیص مدل است؛
    \item $T_M$ هزینه زمانی یا تأخیر مدل است؛
    \item $R_M$ میزان مصرف منابع است؛
    \item و $\alpha$، $\beta$ و $\gamma$ وزن اهمیت هر معیار هستند.
\end{itemize}

این رابطه یک مدل ریاضی دقیق برای انتخاب در پروژه نیست، بلکه
بیان مفهومی این نکته است که کمینه‌کردن \lr{WER} به‌تنهایی هدف کافی
برای یک چت‌بات صوتی محلی نیست.


% ---------------------------------------------------------
\subsection{جمع‌بندی پژوهش‌های چندزبانه و کم‌منبع}

مرور پژوهش‌های این بخش چند نتیجه اصلی را نشان می‌دهد.

نخست، پیش‌آموزش چندزبانه می‌تواند وابستگی به داده برچسب‌خورده
زبان هدف را کاهش دهد و مدل‌هایی مانند XLSR نمونه موفق این
رویکرد هستند. :contentReference[oaicite:0]{index=0}

دوم، حضور یک زبان در یک مدل چندزبانه به‌تنهایی تضمین‌کننده
عملکرد خوب نیست؛ زبان‌هایی که در پیش‌آموزش حضور محدودتری دارند
می‌توانند همچنان به سازگار‌سازی اختصاصی نیاز داشته باشند.
پژوهش‌های انتقال میان‌زبانی نیز نشان داده‌اند که داده زبان‌های
مشابه می‌تواند به این سازگار‌سازی کمک کند. :contentReference[oaicite:1]{index=1}

سوم، مطالعات روی \lr{Whisper} نشان داده‌اند که روش‌هایی مانند
تنظیم اختصاصی، وزن‌دهی زبان و سازگار‌سازی کم‌منبع می‌توانند
عملکرد مدل را بهبود دهند. :contentReference[oaicite:2]{index=2}

در نهایت، پژوهش‌های انجام‌شده روی فارسی و سایر زبان‌های
Perso-Arabic نشان می‌دهند که کیفیت و مرتبط‌بودن داده ممکن است
از اندازه مدل اهمیت بیشتری داشته باشد؛ یک مدل ۳۰۰ میلیون
پارامتری در یک مطالعه اخیر روی فارسی از \lr{Whisper Large-v3}
بهتر عمل کرده است. :contentReference[oaicite:3]{index=3}

این یافته‌ها مبنای تصمیم پروژه حاضر برای مقایسه مستقیم چند
خانواده مدل با اندازه‌ها و روش‌های آموزشی متفاوت فراهم می‌کنند.
% =========================================================
% 3-4
% =========================================================
\section{پژوهش‌های مرتبط با مدل‌های زبانی فارسی}

با گسترش مدل‌های زبانی بزرگ، بخش قابل توجهی از پژوهش‌ها بر
زبان انگلیسی و سایر زبان‌های پرمنبع متمرکز بوده است.
با این حال، عملکرد این مدل‌ها در زبان فارسی الزاماً مشابه
عملکرد آن‌ها در انگلیسی نیست.

تفاوت در حجم داده‌های آموزشی، ساختار نوشتاری زبان فارسی،
کاربرد نیم‌فاصله، شکل‌های مختلف نویسه‌ها، واژگان فرهنگی و
دانش بومی از عواملی هستند که می‌توانند بر عملکرد مدل‌های
زبانی در فارسی اثرگذار باشند.

به همین دلیل، در سال‌های اخیر چند پژوهش با هدف ارزیابی جامع
مدل‌های زبانی روی وظایف فارسی انجام شده‌اند.
همچنین مدل‌هایی مانند
\lr{Dorna}
با هدف سازگارکردن مدل‌های پایه با داده‌های فارسی توسعه یافته‌اند.

در ادامه، مهم‌ترین مطالعات مرتبط با موضوع پروژه حاضر بررسی
می‌شوند.


% ---------------------------------------------------------
\subsection{ارزیابی جامع مدل‌های زبانی در زبان فارسی}

یکی از مطالعات مهم در این حوزه توسط
\lr{Abaskohi et al.}
انجام شده است
\cite{abaskohi2024persianllm}.

این پژوهش یکی از نخستین مطالعات جامع برای بررسی عملکرد
مدل‌های زبانی بزرگ روی مجموعه‌ای متنوع از وظایف فارسی است.
پژوهشگران مدل‌هایی مانند
\lr{GPT-3.5-turbo}،
\lr{GPT-4}
و
\lr{OpenChat-3.5}
را در چند گروه از وظایف مورد ارزیابی قرار دادند.

وظایف مورد بررسی شامل موارد زیر بودند:

\begin{itemize}

    \item وظایف کلاسیک پردازش زبان طبیعی؛

    \item مسائل استدلالی؛

    \item پرسش‌های مبتنی بر دانش عمومی؛

    \item و مسائل ریاضی و آموزشی.

\end{itemize}

از آنجا که مجموعه‌های ارزیابی فارسی برای برخی وظایف استدلالی
محدود بودند، پژوهشگران دو مجموعه جدید شامل پرسش‌های ریاضی
دوره ابتدایی و پرسش‌های آزمون‌های ورودی مدارس ایجاد کردند.

یکی از نتایج مهم این مطالعه آن بود که مدل‌های بزرگ، به‌ویژه
GPT-4، در وظایف نیازمند استدلال و دانش عمومی عملکرد مناسبی
داشتند.
با این حال، در تعدادی از وظایف کلاسیک پردازش زبان طبیعی،
مدل‌های کوچک‌تر که به‌طور اختصاصی برای همان وظیفه تنظیم شده
بودند، عملکرد بهتری ارائه کردند.

همچنین پژوهشگران مشاهده کردند که در برخی موارد، ترجمه ورودی
فارسی به انگلیسی پیش از ارسال آن به GPT-3.5 باعث بهبود عملکرد
می‌شد.

این نتیجه نشان می‌دهد که توانایی یک مدل در زبان فارسی می‌تواند
از توانایی همان مدل در زبان انگلیسی ضعیف‌تر باشد و تنها اندازه
مدل تضمین‌کننده عملکرد مناسب در زبان هدف نیست.

\textbf{ارتباط با پروژه حاضر:}
این یافته مستقیماً با رویکرد پروژه حاضر مرتبط است.
در این پروژه نیز انتخاب مدل زبانی صرفاً بر اساس تعداد
پارامترها انجام نمی‌شود و مدل‌ها به‌صورت مستقیم روی داده‌های
فارسی از نظر کیفیت پاسخ و هزینه اجرایی مقایسه می‌شوند.


% ---------------------------------------------------------
\subsection{نیاز به معیارهای بومی برای فارسی}

یکی از مشکلات مهم ارزیابی مدل‌های زبانی فارسی این است که
ترجمه مستقیم معیارهای انگلیسی لزوماً توانایی واقعی مدل در
زبان فارسی را اندازه‌گیری نمی‌کند.

برای مثال، یک مدل ممکن است در پرسش‌های عمومی ترجمه‌شده عملکرد
مناسبی داشته باشد، اما در موضوعاتی مانند تاریخ ایران،
ادبیات فارسی، اصطلاحات رایج یا دانش فرهنگی عملکرد ضعیف‌تری
نشان دهد.

به همین دلیل، پژوهش‌های جدیدتر به سمت توسعه معیارهایی حرکت
کرده‌اند که به‌طور خاص برای زبان و فرهنگ فارسی طراحی شده‌اند.


% ---------------------------------------------------------
\subsection{معیارهای PeKA و PK-BETS}

در پژوهش
\lr{Hosseinbeigi et al.}
یک چارچوب گسترده‌تر برای ارزیابی مدل‌های زبانی فارسی ارائه
شده است
\cite{hosseinbeigi2025persianllm}.

در این مطالعه دو معیار جدید با نام‌های
\lr{PeKA}
و
\lr{PK-BETS}
برای ارزیابی توانایی مدل‌ها در زمینه‌های مختلف فارسی معرفی
شدند.

این مجموعه‌ها شامل موضوعاتی مانند:

\begin{itemize}

    \item تاریخ و فرهنگ ایران؛

    \item ادبیات فارسی؛

    \item دانش عمومی بومی؛

    \item درک زبان فارسی؛

    \item و وظایفی نیازمند شناخت بافت فرهنگی.

\end{itemize}

هستند.

هدف اصلی این پژوهش آن بود که ارزیابی مدل‌ها تنها به ترجمه
مجموعه‌داده‌های انگلیسی محدود نشود و دانش بومی و زبانی نیز
در نظر گرفته شود.

نتایج این مطالعه نشان می‌دهد که حتی مدل‌های قدرتمند عمومی
هنوز در برخی وظایف نیازمند دانش فرهنگی و زبانی فارسی با
محدودیت مواجه هستند.

\textbf{ارتباط با پروژه حاضر:}
از آنجا که یک چت‌بات فارسی باید پاسخ‌هایی متناسب با زبان
و بافت کاربران فارسی‌زبان تولید کند، ارزیابی مدل تنها بر اساس
معیارهای عمومی کافی نیست.
این موضوع یکی از محدودیت‌های ارزیابی مدل زبانی در پروژه حاضر
نیز به شمار می‌رود و باید در فصل بحث و تحلیل ذکر شود.


% ---------------------------------------------------------
\subsection{معیار FarsEval-PKBETS}

معیار دیگری که برای ارزیابی گسترده مدل‌های فارسی توسعه یافته
است،
\lr{FarsEval-PKBETS}
است
\cite{shamsfard2025farseval}.

این مجموعه شامل حدود ۴۰۰۰ پرسش و پاسخ در قالب‌های مختلف
است.
پرسش‌ها حوزه‌های متنوعی را پوشش می‌دهند، از جمله:

\begin{itemize}

    \item پزشکی؛
    \item حقوق؛
    \item دین؛
    \item زبان فارسی؛
    \item دانش دانشنامه‌ای؛
    \item ترجیحات انسانی؛
    \item دانش اجتماعی؛
    \item اخلاق؛
    \item سوگیری؛
    \item و تولید متن.
\end{itemize}

یکی از ویژگی‌های مهم این مجموعه توجه به ملاحظات زبانی،
فرهنگی و محلی مرتبط با زبان فارسی و ایران است.

در این پژوهش، مدل‌هایی از جمله
\lr{Llama3-70B}،
\lr{PersianMind}
و
\lr{Dorna}
روی این معیار ارزیابی شدند.

نتایج گزارش‌شده نشان دادند که میانگین دقت این مدل‌ها کمتر
از ۵۰ درصد بوده است.
به عبارت دیگر، حتی مدل‌های بزرگ و مدل‌های تنظیم‌شده فارسی
هنوز فاصله قابل توجهی با عملکرد کامل روی مجموعه‌ای گسترده
از وظایف بومی دارند.

\textbf{ارتباط با پروژه حاضر:}
این نتیجه نشان می‌دهد که انتخاب یک مدل فارسی‌محور نیز
به‌تنهایی تضمین‌کننده کیفیت بالا نیست.
بنابراین مدل‌ها باید روی نوع داده و کاربرد واقعی مورد نظر
ارزیابی شوند.


% ---------------------------------------------------------
\subsection{مدل Dorna}

یکی از مدل‌های مرتبط با پروژه حاضر،
خانواده
\lr{Dorna}
است که توسط
\lr{Part AI}
توسعه یافته است
\cite{dorna2modelcard}.

مدل‌های این خانواده از نوع
\lr{Decoder-Only}
هستند و به‌طور ویژه با استفاده از داده‌های فارسی آموزش یا
تنظیم تکمیلی شده‌اند.

نسخه مورد استفاده در پروژه حاضر:

\begin{center}
\lr{PartAI/\lr{Dorna2}-\lr{Llama}3.1-8B-Instruct}
\end{center}

است.

این مدل بر پایه
\lr{Meta Llama 3.1 8B Instruct}
ساخته شده و سپس برای زبان فارسی تنظیم شده است.

طبق مستندات رسمی مدل، \lr{Dorna2} به‌عنوان یک مدل ۸ میلیارد
پارامتری برای کاربردهای گفت‌وگویی فارسی ارائه شده است.

رویکرد مورد استفاده در Dorna نمونه‌ای از
\textbf{سازگارکردن یک مدل پایه عمومی با زبان هدف}
است.

به‌جای آموزش یک مدل بزرگ از ابتدا، می‌توان از دانش عمومی
یک مدل پایه استفاده کرده و سپس با داده‌های زبان فارسی رفتار
آن را برای کاربرد فارسی بهبود داد.

این رویکرد از نظر هزینه آموزش نیز می‌تواند عملی‌تر باشد.


% ---------------------------------------------------------
\subsection{\lr{Dorna2} در مقایسه با مدل پایه}

ارزیابی‌های ارائه‌شده در مستندات \lr{Dorna2} نشان می‌دهند که
نسخه تنظیم‌شده فارسی در تعدادی از معیارها نسبت به
\lr{Llama 3.1 8B Instruct}
پایه عملکرد بهتری دارد.

این موضوع نشان می‌دهد که افزودن داده و تنظیم اختصاصی فارسی
می‌تواند در برخی وظایف باعث بهبود مدل شود.

با این حال، تفاوت عملکرد در همه معیارها یکسان نیست.
بنابراین نمی‌توان نتیجه گرفت که تنظیم فارسی در تمام کاربردها
همیشه بهبود ایجاد می‌کند.

\textbf{ارتباط با پروژه حاضر:}
این موضوع یکی از دلایل انتخاب \lr{Dorna2} به‌عنوان یکی از مدل‌های
بنچمارک پروژه حاضر بوده است.
با مقایسه آن با مدل‌هایی مانند Gemma و Qwen می‌توان بررسی کرد
که آیا تخصص بیشتر در زبان فارسی می‌تواند در مقابل اختلاف
اندازه مدل و هزینه محاسباتی مزیت ایجاد کند.


% ---------------------------------------------------------
\subsection{مدل‌های عمومی چندزبانه در برابر مدل‌های فارسی‌محور}

در طراحی سامانه فارسی دو راهبرد کلی را می‌توان در نظر گرفت.

راهبرد اول استفاده از یک مدل عمومی چندزبانه است.
مدل‌هایی مانند
\lr{Gemma}،
\lr{Qwen}
و
\lr{Llama}
با حجم بزرگی از داده‌ها و چندین زبان آموزش دیده‌اند و دانش
عمومی گسترده‌ای دارند.

راهبرد دوم استفاده از مدلی است که به‌طور اختصاصی با داده‌های
فارسی تنظیم شده باشد.

مزیت مدل عمومی می‌تواند دانش گسترده‌تر و توانایی بهتر در
وظایف عمومی باشد.
در مقابل، مدل فارسی‌محور ممکن است در درک ساختار زبان،
واژگان و کاربردهای فارسی عملکرد بهتری ارائه دهد.

اما این رابطه قطعی نیست.

نتایج پژوهش‌های ارزیابی فارسی نشان داده‌اند که مدل‌های بزرگ
عمومی در برخی وظایف بهتر عمل می‌کنند، در حالی که در برخی
وظایف مدل‌های کوچک‌تر یا تنظیم‌شده اختصاصی می‌توانند عملکرد
رقابتی یا حتی بهتر داشته باشند
\cite{abaskohi2024persianllm}.

در نتیجه، انتخاب میان این دو گروه باید مبتنی بر ارزیابی
تجربی باشد.


% ---------------------------------------------------------
\subsection{ارزیابی فرهنگی مدل‌های فارسی}

یکی از جهت‌گیری‌های جدید پژوهش در این حوزه، بررسی توانایی
مدل‌ها در درک فرهنگ و دانش بومی ایران است.

معیار
\lr{MELAC}
با هدف ارزیابی گسترده مدل‌های زبانی از نظر سازگاری فرهنگی
در زبان فارسی ارائه شده است
\cite{farsi2025melac}.

در این پژوهش ۱۹ مجموعه ارزیابی جدید در موضوعاتی مانند:

\begin{itemize}

    \item قوانین ایران؛
    \item دستور زبان فارسی؛
    \item اصطلاحات فارسی؛
    \item دانش فرهنگی؛
    \item و آزمون‌های ورودی دانشگاهی
\end{itemize}

ایجاد شدند.

در مجموع ۴۱ مدل زبانی بزرگ روی این مجموعه‌ها ارزیابی شدند.

این نوع پژوهش‌ها نشان می‌دهند که کیفیت یک مدل زبانی نباید
تنها بر اساس روان‌بودن متن خروجی ارزیابی شود.
مدل ممکن است پاسخ بسیار روانی تولید کند اما دانش فرهنگی یا
اطلاعات بومی نادرستی ارائه دهد.


% ---------------------------------------------------------
\subsection{ایمنی و هم‌راستایی در زبان فارسی}

موضوع دیگری که در پژوهش‌های جدید مورد توجه قرار گرفته،
\textbf{هم‌راستایی}
\LTRfootnote{Alignment}
مدل‌های زبانی فارسی است.

در پژوهش
\lr{ELAB}
یک چارچوب ارزیابی برای بررسی ابعاد مختلف هم‌راستایی مدل‌های
فارسی ارائه شده است
\cite{pourbahman2025elab}.

این چارچوب حوزه‌هایی مانند:

\begin{itemize}

    \item ایمنی؛
    \item عدالت و سوگیری؛
    \item هنجارهای اجتماعی؛
    \item و محتوای آسیب‌زا
\end{itemize}

را مورد بررسی قرار می‌دهد.

برای این هدف، علاوه بر ترجمه برخی معیارهای انگلیسی، مجموعه‌های
جدید فارسی نیز ایجاد شده‌اند تا مسائل فرهنگی و اجتماعی بومی
بهتر پوشش داده شوند.

اگرچه ایمنی و هم‌راستایی هدف اصلی پروژه حاضر نیستند،
این مطالعات نشان می‌دهند که ارزیابی یک چت‌بات واقعی باید
در آینده فراتر از کیفیت زبانی و سرعت پاسخ باشد.


% ---------------------------------------------------------
\subsection{چالش ارزیابی کیفیت پاسخ آزاد}

بر خلاف وظایفی مانند دسته‌بندی متن، ارزیابی خروجی یک مدل
گفت‌وگویی دشوارتر است؛ زیرا ممکن است برای یک پرسش چند پاسخ
صحیح وجود داشته باشد.

در چنین شرایطی، تطابق دقیق پاسخ مدل با یک متن مرجع می‌تواند
معیار مناسبی نباشد.

برای مثال، دو پاسخ:

\begin{center}
«تهران پایتخت ایران است.»
\end{center}

و:

\begin{center}
«پایتخت کشور ایران شهر تهران است.»
\end{center}

از نظر معنایی تقریباً یکسان هستند، اما تطابق واژه‌به‌واژه آن‌ها
کامل نیست.

به همین دلیل، در ارزیابی مدل‌های زبانی می‌توان از معیارهای
معنایی مبتنی بر بازنمایی برداری
\LTRfootnote{Embedding-based Similarity}
در کنار معیارهای واژگانی مانند
\lr{ROUGE}
استفاده کرد.

در پروژه حاضر، دقت برای پرسش‌های چندگزینه‌ای و ترکیب تطابق دقیق،
\lr{F1} توکنی و شباهت معنایی برای درک مطلب استفاده شده است.
شباهت معنایی و \lr{ROUGE-L} نیز در آزمایش مقدماتی تولید آزاد
به کار رفتند.

با این حال، هیچ‌یک از این معیارها به‌تنهایی کیفیت کامل یک
پاسخ گفت‌وگویی را نشان نمی‌دهند.
در یک ارزیابی جامع‌تر، بررسی انسانی یا داوری مبتنی بر مجموعه‌ای
از معیارهای صحت، مرتبط‌بودن، روانی و مفیدبودن می‌تواند
اطلاعات تکمیلی ارائه کند.


% ---------------------------------------------------------
\subsection{مدل زبانی در شرایط اجرای محلی}

بخش بزرگی از پژوهش‌های ارزیابی مدل‌های زبانی بر کیفیت خروجی
تمرکز دارند.
اما برای یک چت‌بات محلی، محدودیت‌های سخت‌افزاری نیز اهمیت
اساسی دارند.

مدلی با کیفیت کمی بهتر ممکن است چند برابر حافظه بیشتری
مصرف کند یا زمان پاسخ آن به اندازه‌ای زیاد باشد که برای
تعامل صوتی مناسب نباشد.

بنابراین در کاربرد مورد نظر این پروژه باید حداقل سه بعد
به‌صورت هم‌زمان بررسی شوند:

\begin{equation}
Q_{\\mathrm{LLM}}
=
f(
Q_{\mathrm{semantic}},
T_{\mathrm{response}},
M_{\mathrm{memory}}
)
\end{equation}

که در آن:

\begin{itemize}

    \item $Q_{\mathrm{semantic}}$ کیفیت معنایی پاسخ؛
    \item $T_{\mathrm{response}}$ زمان پاسخ؛
    \item و $M_{\mathrm{memory}}$ میزان حافظه مصرف‌شده است.
\end{itemize}

در عمل معیارهایی مانند سرعت تولید توکن نیز باید در کنار این
موارد اندازه‌گیری شوند.

این دیدگاه با نتایج پژوهش‌های فارسی نیز سازگار است؛ زیرا
مدل بزرگ‌تر همیشه بهترین عملکرد را برای همه وظایف ارائه
نمی‌کند.


% ---------------------------------------------------------
\subsection{جمع‌بندی پژوهش‌های مدل زبانی فارسی}

مرور مطالعات مرتبط چند نتیجه مهم را نشان می‌دهد.

نخست، عملکرد مدل‌های زبانی در فارسی هنوز با عملکرد آن‌ها
در زبان‌های پرمنبع کاملاً برابر نیست و در برخی مطالعات
استفاده از نسخه انگلیسی ورودی باعث بهبود عملکرد مدل شده است
\cite{abaskohi2024persianllm}.

دوم، ارزیابی فارسی باید علاوه بر معیارهای ترجمه‌شده عمومی،
دانش بومی، فرهنگ، ادبیات، هنجارهای اجتماعی و ویژگی‌های خاص
زبان فارسی را نیز پوشش دهد.

سوم، مدل‌های فارسی‌محور مانند Dorna می‌توانند از طریق تنظیم
مدل پایه روی داده‌های فارسی برخی توانایی‌های مدل را بهبود
دهند، اما این مزیت در همه وظایف تضمین‌شده نیست.

در نهایت، برای کاربرد مورد نظر این پروژه، کیفیت پاسخ تنها
یکی از معیارهای انتخاب مدل است.
مدل منتخب باید بتواند با تأخیر مناسب و مصرف حافظه قابل قبول
در کنار \lr{ASR} و \lr{TTS} روی سخت‌افزار محلی اجرا شود.

بر همین اساس، در پروژه حاضر مدل‌هایی با اندازه‌ها و رویکردهای
متفاوت به‌صورت مستقیم روی یک مجموعه فارسی مشترک ارزیابی
شده‌اند.
% =========================================================
% 3-5
% =========================================================
\section{مدل‌های زبانی سبک و اجرای محلی}

افزایش اندازه مدل‌های زبانی در سال‌های اخیر باعث بهبود توانایی
آن‌ها در بسیاری از وظایف شده است، اما این افزایش اندازه هزینه
محاسباتی قابل توجهی نیز ایجاد می‌کند.

مدل‌هایی با ده‌ها یا صدها میلیارد پارامتر معمولاً برای اجرا به
حافظه زیاد، پردازنده‌های گرافیکی قدرتمند و زیرساخت محاسباتی
مناسب نیاز دارند.
این محدودیت باعث شده است استفاده از چنین مدل‌هایی در
سامانه‌های محلی یا دستگاه‌های شخصی دشوار باشد.

در مقابل، مدل‌های زبانی کوچک‌تر می‌توانند امکان اجرای محلی،
کاهش هزینه و کاهش وابستگی به سرویس‌های ابری را فراهم کنند.

در کاربرد مورد نظر این پروژه، این موضوع اهمیت بیشتری دارد؛
زیرا مدل زبانی تنها مؤلفه سامانه نیست و باید در کنار مدل
تشخیص گفتار و تبدیل متن به گفتار روی یک سیستم مشترک اجرا شود.


% ---------------------------------------------------------
\subsection{انگیزه استفاده از مدل‌های کوچک‌تر}

یک مدل بزرگ‌تر معمولاً ظرفیت بیشتری برای ذخیره و پردازش
الگوهای زبانی دارد، اما افزایش تعداد پارامترها مستقیماً باعث
افزایش حافظه مورد نیاز برای ذخیره وزن‌ها می‌شود.

اگر تعداد پارامترها را با $N$ و تعداد بیت مورد استفاده برای
هر پارامتر را با $b$ نمایش دهیم، حافظه خام مورد نیاز برای
ذخیره وزن‌ها تقریباً برابر است با:

\begin{equation}
M_{\mathrm{weights}}
\approx
\frac{N \times b}{8}
\end{equation}

برای مثال، اگر وزن‌ها با دقت 16 بیت ذخیره شوند، یک مدل
8 میلیارد پارامتری تنها برای وزن‌های خود تقریباً به:

\begin{equation}
8 \times 10^9
\times
2
\approx
16 \; \mathrm{GB}
\end{equation}

حافظه نیاز خواهد داشت.

این مقدار شامل حافظه مورد نیاز برای کش توجه، فعال‌سازی‌ها،
بافرهای محاسباتی و سایر اجزای زمان اجرا نیست.

بنابراین، استفاده از مدلی با پارامتر کمتر می‌تواند امکان اجرای
مدل روی سخت‌افزارهای معمولی‌تر را فراهم کند.


% ---------------------------------------------------------
\subsection{اجرای مدل روی دستگاه محلی}

اجرای مستقیم مدل زبانی روی دستگاه کاربر با عنوان‌هایی مانند

\textbf{اجرای محلی}
\LTRfootnote{Local Inference}

یا

\textbf{استنتاج روی دستگاه}
\LTRfootnote{On-device Inference}

شناخته می‌شود.

در این حالت، درخواست کاربر برای تولید پاسخ به یک سرویس خارجی
ارسال نمی‌شود و محاسبات اصلی روی سخت‌افزار محلی انجام می‌شوند.

این رویکرد چند مزیت بالقوه دارد:

\begin{itemize}

    \item کاهش وابستگی به اتصال اینترنت؛

    \item امکان استفاده آفلاین؛

    \item کنترل بیشتر بر زمان پاسخ؛

    \item کاهش وابستگی به سرویس‌های خارجی؛

    \item و حفظ محلی داده‌های ورودی و تاریخچه مکالمه.

\end{itemize}

با این حال، محدودیت اصلی این روش توان محاسباتی و حافظه
دستگاه است.
به همین دلیل، توسعه روش‌هایی برای کاهش اندازه مدل و هزینه
استنتاج به یکی از موضوعات مهم پژوهشی تبدیل شده است.


% ---------------------------------------------------------
\subsection{کوانتیزه‌سازی پس از آموزش}

یکی از مهم‌ترین روش‌های کاهش هزینه اجرای مدل،
\textbf{کوانتیزه‌سازی پس از آموزش}
\LTRfootnote{Post-Training Quantization (PTQ)}
است.

در این روش، پس از پایان آموزش مدل، وزن‌های ذخیره‌شده با دقت
عددی پایین‌تری نمایش داده می‌شوند.

برای مثال:

\begin{center}
FP16
$\longrightarrow$
INT8
$\longrightarrow$
INT4
\end{center}

کاهش دقت از 16 بیت به 4 بیت از نظر نظری می‌تواند حافظه لازم
برای ذخیره وزن‌ها را تقریباً به یک‌چهارم کاهش دهد.

مطابق رابطه حافظه وزن‌ها که پیش‌تر ارائه شد، وزن‌های یک مدل
\lr{9B} در دقت \lr{FP16} تقریباً
\lr{18 GB} فضا نیاز دارند، در حالی که مقدار نظری آن در دقت 4 بیتی
حدود \lr{4.5 GB} است. برای یک مدل \lr{2B} نیز این مقادیر به‌ترتیب
حدود \lr{4 GB} و \lr{1 GB} هستند. این محاسبه فقط وزن‌ها را در نظر
می‌گیرد؛ حافظه واقعی بیشتر است، زیرا مقیاس‌های کوانتیزه‌سازی،
فعال‌سازی‌ها، فضای کاری کرنل‌ها و حافظه نهان \lr{KV} نیز باید در
حافظه قرار گیرند.

با این حال، تبدیل ساده وزن‌ها به تعداد بیت کمتر می‌تواند
باعث کاهش دقت مدل شود.
به همین دلیل، روش‌های پیشرفته‌تری برای تعیین نحوه کوانتیزه‌سازی
وزن‌ها توسعه یافته‌اند.


% ---------------------------------------------------------
\subsection{کوانتیزه‌سازی چهاربیتی \lr{NF4}}

در ارزیابی‌های اصلی این پروژه از قالب چهاربیتی
\lr{NormalFloat 4 (NF4)} استفاده شد. این قالب در پژوهش \lr{QLoRA}
برای وزن‌هایی با توزیع تقریباً نرمال معرفی شده است و سطوح نمایش
را متناسب با توزیع آماری وزن‌ها تعیین می‌کند
\cite{dettmers2023qlora}. هدف آن این است که در بودجه ثابت 4 بیت،
خطای بازنمایی وزن‌ها نسبت به یک قالب یکنواخت کاهش یابد.

در این شیوه باید میان
\textbf{دقت ذخیره‌سازی}
و
\textbf{دقت محاسبات}
تمایز قائل شد. وزن‌ها به‌صورت 4 بیتی در حافظه نگهداری می‌شوند،
اما هنگام اجرای ضرب‌های ماتریسی، بلوک مورد نیاز به یک دقت محاسباتی
بالاتر بازگردانی می‌شود. بنابراین عبارت «مدل 4 بیتی» لزوماً به این
معنا نیست که تمام فعال‌سازی‌ها، حافظه نهان \lr{KV} و عملیات عددی
سامانه نیز با دقت 4 بیت انجام می‌شوند.

تکنیک
\textbf{کوانتیزه‌سازی دوگانه}
\LTRfootnote{Double Quantization}
نیز می‌تواند خود ضرایب مقیاس کوانتیزه‌سازی را فشرده کند و سربار
حافظه را بیشتر کاهش دهد \cite{dettmers2023qlora}. با این حال، استفاده
از این قابلیت به تنظیمات کتابخانه و پیاده‌سازی وابسته است و نباید
صرفاً از ذکر \lr{NF4} نتیجه گرفت که کوانتیزه‌سازی دوگانه نیز فعال
بوده است.

انتخاب \lr{NF4} در این پروژه یک تصمیم استقراری بود: همه مدل‌ها باید
در شرایط یکسان روی یک \lr{GPU} محلی قابل بارگذاری و ارزیابی می‌بودند.
کوانتیزه‌سازی چهاربیتی امکان قرارگرفتن مدل‌های \lr{2B} تا \lr{9B}
در محدوده حافظه سخت‌افزار محلی را فراهم کرد و در عین حال، مقایسه را
از اثر متفاوت‌بودن دقت عددی میان مدل‌ها مصون نگه داشت.


% ---------------------------------------------------------
\subsection{روش \lr{GPTQ}}

یکی از روش‌های شناخته‌شده برای کوانتیزه‌سازی مدل‌های زبانی،
\lr{GPTQ}
است که توسط
\lr{Frantar et al.}
معرفی شده است
\cite{frantar2022gptq}.

هدف \lr{GPTQ} کاهش دقت عددی وزن‌های مدل با کمترین افت ممکن در
کیفیت خروجی است.

این روش از اطلاعات تقریبی مرتبه دوم برای تعیین خطای ناشی از
کوانتیزه‌کردن وزن‌ها استفاده می‌کند.

پژوهشگران نشان دادند که مدل‌های بسیار بزرگ را می‌توان با
وزن‌های 3 یا 4 بیتی فشرده کرد، در حالی که افت عملکرد نسبت
به نسخه اصلی محدود باقی بماند.

در همان مطالعه، کوانتیزه‌سازی مدل‌های بزرگ موجب کاهش چشمگیر
حافظه و همچنین افزایش سرعت استنتاج در سخت‌افزارهای مورد
آزمایش شد.

\textbf{ارتباط با پروژه حاضر:}
استفاده از تکنیک‌هایی مانند \lr{GPTQ} می‌تواند امکان اجرای مدلی
را فراهم کند که در حالت FP16 از ظرفیت حافظه سیستم خارج باشد.
با این حال، اثر کوانتیزه‌سازی بر کیفیت فارسی باید به‌صورت
تجربی بررسی شود و نمی‌توان فرض کرد که افت کیفیت برای تمام
مدل‌ها و زبان‌ها یکسان است.


% ---------------------------------------------------------
\subsection{روش \lr{AWQ}}

رویکرد دیگری با عنوان
\lr{Activation-aware Weight Quantization (AWQ)}
توسط
\lr{Lin et al.}
ارائه شده است
\cite{lin2023awq}.

ایده اصلی این روش آن است که تمام وزن‌های مدل از نظر تأثیر
بر خروجی اهمیت یکسانی ندارند.

پژوهشگران نشان دادند که بخش کوچکی از کانال‌های وزن که با
فعال‌سازی‌های بزرگ مرتبط هستند، نقش بیشتری در حفظ عملکرد
مدل دارند.

در \lr{AWQ} به‌جای استفاده از دقت ترکیبی پیچیده، مقیاس برخی
کانال‌های مهم تغییر داده می‌شود تا خطای کوانتیزه‌سازی آن‌ها
کاهش پیدا کند.

مزیت مهم این روش آن است که برای تعیین وزن‌های مهم از
آمار فعال‌سازی استفاده می‌شود و نیاز به آموزش مجدد کامل
مدل ندارد.

در نتایج ارائه‌شده، \lr{AWQ} امکان استنتاج مدل‌های 4 بیتی روی
سخت‌افزارهای رومیزی و حتی برخی دستگاه‌های محلی را با افزایش
سرعت قابل توجه فراهم کرده است
\cite{lin2023awq}.

این پژوهش نشان می‌دهد که اجرای محلی مدل‌های زبانی تنها به
کوچک‌کردن تعداد پارامترها محدود نیست و نحوه نمایش وزن‌ها نیز
نقش مهمی در امکان استقرار مدل دارد.


% ---------------------------------------------------------
\subsection{توازن میان حافظه و کیفیت}

کوانتیزه‌سازی همواره یک مسئله توازن است.

کاهش تعداد بیت‌ها باعث کاهش حافظه و پهنای باند مورد نیاز
برای انتقال وزن‌ها می‌شود، اما می‌تواند خطای عددی بیشتری
ایجاد کند.

اثرهای اصلی کوانتیزه‌سازی را می‌توان در چهار محور خلاصه کرد:

\begin{itemize}

    \item
    \textbf{حافظه:}
    کاهش دقت وزن‌ها بیشترین اثر مستقیم را بر حافظه وزن‌ها دارد،
    اما مصرف کل حافظه به طول ورودی، اندازه دسته، حافظه نهان \lr{KV}
    و سربار چارچوب اجرا نیز وابسته است.

    \item
    \textbf{کیفیت:}
    نگاشت وزن‌های پیوسته به تعداد محدودی سطح گسسته خطای تقریب ایجاد
    می‌کند. این خطا ممکن است بر مدل‌ها، لایه‌ها، زبان‌ها و وظایف
    مختلف اثر یکسانی نداشته باشد؛ ازاین‌رو کیفیت نسخه کوانتیزه‌شده
    باید روی وظایف هدف دوباره اندازه‌گیری شود.

    \item
    \textbf{سرعت:}
    کاهش انتقال داده از حافظه می‌تواند استنتاج را سریع‌تر کند،
    به‌ویژه هنگامی که گلوگاه پهنای باند حافظه باشد. در مقابل، هزینه
    بازکوانی وزن‌ها و نبود کرنل بهینه روی برخی سخت‌افزارها می‌تواند
    این مزیت را کاهش دهد؛ بنابراین 4 بیتی‌بودن به‌تنهایی تضمین‌کننده
    افزایش سرعت نیست.

    \item
    \textbf{امکان استقرار:}
    مهم‌ترین اثر عملی در این پروژه، تبدیل «مدلی خارج از ظرفیت حافظه»
    به «مدلی قابل بارگذاری روی \lr{GPU} محلی» بود. این مزیت امکان
    مقایسه مدل‌های بزرگ‌تر را بدون استفاده از سرویس ابری فراهم کرد.

\end{itemize}

به‌صورت مفهومی می‌توان این مسئله را چنین نمایش داد:

\begin{equation}
Q_{\mathrm{deployment}}
=
f(
M_{\mathrm{memory}},
S_{\mathrm{speed}},
A_{\mathrm{quality}}
)
\end{equation}

که در آن:

\begin{itemize}

    \item $M_{\mathrm{memory}}$ حافظه مورد نیاز؛

    \item $S_{\mathrm{speed}}$ سرعت استنتاج؛

    \item و $A_{\mathrm{quality}}$ کیفیت مدل است.

\end{itemize}

یک روش فشرده‌سازی مناسب باید حافظه و زمان اجرا را کاهش دهد
بدون آنکه کیفیت مدل به میزان غیرقابل قبول افت کند.

پژوهش‌هایی مانند \lr{GPTQ} و \lr{AWQ} نشان داده‌اند که کوانتیزه‌سازی
4 بیتی در بسیاری از مدل‌ها می‌تواند نقطه توازن مناسبی میان
این سه عامل ایجاد کند
\cite{frantar2022gptq,lin2023awq}.


% ---------------------------------------------------------
\subsection{مرحله Prefill و مرحله Decode}

زمان استنتاج یک مدل زبانی را می‌توان به دو مرحله اصلی تقسیم
کرد.

مرحله نخست،
\textbf{Prefill}
است.
در این مرحله تمام توکن‌های ورودی کاربر و تاریخچه مکالمه
پردازش می‌شوند.

مرحله دوم،
\textbf{Decode}
است که در آن مدل توکن‌های خروجی را به‌صورت خودبازگشتی
تولید می‌کند.

به‌طور مفهومی:

\begin{center}
Prompt
$\longrightarrow$
Prefill
$\longrightarrow$
Token 1
$\longrightarrow$
Token 2
$\longrightarrow$
$\cdots$
\end{center}

هزینه این دو مرحله رفتار یکسانی ندارد.
مرحله Prefill معمولاً شامل محاسبات ماتریسی بزرگ روی کل
دنباله ورودی است، در حالی که مرحله Decode در هر گام تنها
یک یا تعداد کمی توکن جدید تولید می‌کند.

به همین دلیل، یک روش بهینه‌سازی ممکن است سرعت Prefill را
افزایش دهد ولی تأثیر متفاوتی بر Decode داشته باشد.

در یک چت‌بات، سرعت Decode اهمیت زیادی دارد زیرا مستقیماً
بر سرعت ظاهرشدن پاسخ برای کاربر اثر می‌گذارد.


% ---------------------------------------------------------
\subsection{سرعت تولید توکن}

یکی از معیارهای متداول برای سنجش سرعت مدل زبانی،
تعداد توکن تولیدشده در ثانیه است:

\begin{equation}
TPS
=
\frac{N_{\mathrm{tokens}}}
     {T_{\mathrm{decode}}}
\end{equation}

که در آن $N_{\mathrm{tokens}}$ تعداد توکن‌های تولیدشده
و $T_{\mathrm{decode}}$ زمان تولید آن‌ها است.

مقدار بیشتر این معیار نشان‌دهنده سرعت بالاتر تولید پاسخ است.

با این حال، مقایسه مستقیم این مقدار بین مطالعات مختلف باید
با احتیاط انجام شود، زیرا عواملی مانند \lr{GPU}، نوع کوانتیزه‌سازی،
طول ورودی، اندازه batch، چارچوب استنتاج و طول پاسخ می‌توانند
بر آن اثر بگذارند.

بنابراین در یک بنچمارک منصفانه بهتر است مدل‌های مختلف روی
سخت‌افزار و تنظیمات یکسان اجرا شوند.


% ---------------------------------------------------------
\subsection{زمان تولید نخستین توکن}

در یک سامانه مکالمه‌ای، یکی دیگر از معیارهای مهم
\textbf{زمان تا نخستین توکن}
\LTRfootnote{\lr{Time to First Token} (\lr{TTFT})}
است.

این معیار فاصله زمانی میان ارسال ورودی به مدل و تولید نخستین
توکن پاسخ را اندازه‌گیری می‌کند.

در حالی که تعداد توکن در ثانیه سرعت ادامه تولید پاسخ را نشان
می‌دهد، \lr{TTFT} مشخص می‌کند کاربر چه مدت منتظر شروع پاسخ می‌ماند.

به‌صورت مفهومی:

\begin{equation}
T_{\mathrm{response}}
=
T_{\\mathrm{TTFT}}
+
T_{\mathrm{decode}}
\end{equation}

در کاربرد تعاملی، یک مدل ممکن است سرعت تولید توکن بالایی داشته
باشد اما زمان آغاز پاسخ آن زیاد باشد.
بنابراین اندازه‌گیری هر دو نوع تأخیر می‌تواند تصویر دقیق‌تری
از تجربه کاربر ارائه کند.


% ---------------------------------------------------------
\subsection{کش کلید و مقدار}

در مدل‌های مبتنی بر \lr{Transformer}، هنگام تولید خودبازگشتی،
محاسبه مجدد تمام اطلاعات مربوط به توکن‌های قبلی در هر مرحله
هزینه زیادی ایجاد می‌کند.

برای کاهش این هزینه معمولاً از
\textbf{KV Cache}
یا کش کلید و مقدار استفاده می‌شود.

در این روش، بردارهای Key و Value مربوط به توکن‌های قبلی
ذخیره می‌شوند و در مراحل بعدی دوباره مورد استفاده قرار
می‌گیرند.

این روش سرعت تولید را افزایش می‌دهد، اما حافظه بیشتری مصرف
می‌کند.

هرچه طول تاریخچه مکالمه افزایش یابد، اندازه KV Cache نیز
بیشتر خواهد شد.

بنابراین در یک چت‌بات چندمرحله‌ای، طول بافت نه‌تنها بر هزینه
محاسباتی بلکه بر مصرف حافظه نیز اثر دارد.


% ---------------------------------------------------------
\subsection{مدل‌های کوچک و تخصصی}

یکی از جهت‌گیری‌های پژوهشی جدید، استفاده از مدل‌های کوچک‌تر
برای وظایف مشخص است.

در چنین رویکردی، هدف لزوماً رقابت با یک مدل بسیار بزرگ در
تمام وظایف عمومی نیست.
در عوض، یک مدل کوچک‌تر برای مجموعه‌ای محدود از وظایف یا
یک دامنه مشخص انتخاب و در صورت نیاز تنظیم می‌شود.

پژوهش‌های جدید در زمینه استقرار محلی نشان می‌دهند که مدل‌های
زیر چند میلیارد پارامتر می‌توانند در برخی وظایف ساختاریافته
عملکرد قابل قبول ارائه کنند، به‌ویژه هنگامی که انتخاب مدل و
تنظیم آن متناسب با کاربرد انجام شود.

این رویکرد برای سامانه‌هایی با محدودیت منابع اهمیت زیادی دارد،
زیرا کاهش اندازه مدل می‌تواند امکان اجرای هم‌زمان چند مؤلفه
هوش مصنوعی را روی یک دستگاه فراهم کند.


% ---------------------------------------------------------
\subsection{مدل‌های سبک در پروژه حاضر}

در پروژه حاضر، مدل زبانی باید هم‌زمان چند شرط را برآورده کند:

\begin{itemize}

    \item توانایی تولید پاسخ فارسی مناسب؛

    \item زمان پاسخ قابل قبول؛

    \item سرعت مناسب تولید توکن؛

    \item قرارگرفتن وزن‌ها و سربار استنتاج در ظرفیت حافظه \lr{GPU} محلی؛

    \item و امکان اجرای هم‌زمان با \lr{ASR} و \lr{TTS}.

\end{itemize}

قابلیت بارگذاری و استنتاج کامل روی یک \lr{GPU} محلی، شرط اولیه
ورود مدل به آزمایش بود. به همین دلیل، به‌جای مدل‌های بسیار بزرگ
که به چند شتاب‌دهنده یا انتقال بخشی از محاسبات به حافظه اصلی نیاز
دارند، مدل‌هایی در محدوده 2 تا 9 میلیارد پارامتر انتخاب شدند.
این مدل‌ها با کوانتیزه‌سازی چهاربیتی \lr{NF4} در محدوده حافظه
سخت‌افزار محلی پروژه قرار می‌گرفتند و می‌شد آن‌ها را تحت یک تنظیم
عددی مشترک مقایسه کرد.

مدل‌های مورد بررسی شامل:

\begin{itemize}

    \item \lr{Gemma 2 2B};

    \item \lr{Gemma 2 9B};

    \item \lr{Qwen2.5 3B};

    \item \lr{Qwen2.5 7B};

    \item و \lr{Dorna2-Llama3.1 8B}

\end{itemize}

هستند.

وجود مدل‌هایی با اندازه‌های متفاوت امکان بررسی مستقیم این
پرسش را فراهم می‌کند که آیا افزایش تعداد پارامترها واقعاً
به بهبود کافی در کیفیت فارسی منجر می‌شود تا افزایش مصرف حافظه
و تأخیر را توجیه کند یا خیر.

بنابراین فهرست مدل‌ها صرفاً بر اساس محبوبیت یا کیفیت گزارش‌شده
آن‌ها تعیین نشد؛ امکان استقرار محلی، پشتیبانی از زبان فارسی، وجود
نسخه دستورمحور و پوشش چند نقطه متفاوت از توازن اندازه و کیفیت،
معیارهای اصلی انتخاب بودند.


% ---------------------------------------------------------
\subsection{انتخاب مدل به‌عنوان مسئله چندمعیاره}

برای یک سامانه محلی، انتخاب مدل مناسب را می‌توان به‌عنوان
یک مسئله چندمعیاره در نظر گرفت.

فرض شود برای مدل $M$ معیارهای زیر اندازه‌گیری شوند:

\begin{equation}
A(M) = \mbox{\rl{کیفیت پاسخ}}
\end{equation}

\begin{equation}
L(M) = \mbox{\rl{زمان تأخیر}}
\end{equation}

\begin{equation}
V(M) = \mbox{\rl{مصرف حافظه}}
\end{equation}

\begin{equation}
S(M) = \mbox{\rl{سرعت تولید توکن}}
\end{equation}

پیش از مقایسه این معیارها، مدل باید قید امکان استقرار را برآورده کند:

\begin{equation}
V_{\mathrm{peak}}(M) + V_{\mathrm{ASR}} + V_{\mathrm{TTS}}
\leq V_{\mathrm{GPU}}
\end{equation}

که در آن $V_{\mathrm{peak}}(M)$ بیشینه حافظه مورد نیاز مدل زبانی و
$V_{\mathrm{GPU}}$ ظرفیت قابل استفاده شتاب‌دهنده محلی است. دو جمله
دیگر سهم مؤلفه‌های تشخیص گفتار و تبدیل متن به گفتار را نشان می‌دهند.
اگر این قید برقرار نباشد، حتی مدلی با کیفیت بالاتر نیز برای اجرای
یکپارچه سامانه گزینه عملی محسوب نمی‌شود.

در این حالت، مدلی که تنها بیشترین کیفیت را دارد الزاماً
بهترین انتخاب عملی نیست.

برای مثال، اگر افزایش اندکی در کیفیت نیازمند دو یا سه برابر
حافظه بیشتر باشد، مدل کوچک‌تر ممکن است برای کاربرد محلی
انتخاب مناسب‌تری باشد.

این مفهوم را می‌توان با مفهوم
\textbf{مرز پارتو}
\LTRfootnote{Pareto Frontier}
توضیح داد.
یک مدل روی مرز پارتو قرار می‌گیرد اگر نتوان یک معیار آن را
بدون بدترشدن حداقل یکی از معیارهای دیگر بهبود داد.

این دیدگاه برای تحلیل نتایج مدل‌های پروژه حاضر مفید است؛
زیرا هدف اصلی دستیابی به بهترین توازن میان کیفیت و کارایی
است، نه صرفاً بیشینه‌کردن یک معیار منفرد.


% ---------------------------------------------------------
\subsection{جمع‌بندی مدل‌های سبک و اجرای محلی}

مطالعات مرتبط با استقرار مدل‌های زبانی نشان می‌دهند که یکی
از مهم‌ترین موانع اجرای محلی، حجم وزن‌ها و هزینه استنتاج است.

روش‌های کوانتیزه‌سازی مانند \lr{GPTQ} امکان کاهش وزن‌ها به 3 یا
4 بیت را با افت محدود عملکرد فراهم کرده‌اند
\cite{frantar2022gptq}.

روش \lr{AWQ} نیز با استفاده از اطلاعات فعال‌سازی تلاش می‌کند
وزن‌های مهم‌تر را هنگام کوانتیزه‌سازی بهتر حفظ کند و برای
استنتاج 4 بیتی روی دستگاه‌های محلی طراحی شده است
\cite{lin2023awq}.

این مطالعات نشان می‌دهند که امکان اجرای یک مدل به تعداد
پارامترهای آن محدود نمی‌شود و نوع نمایش وزن‌ها، طول بافت،
KV Cache، سخت‌افزار و چارچوب استنتاج نیز عوامل مهمی هستند.

در نهایت، برای کاربرد مورد نظر پروژه حاضر، مدل مناسب باید
در یک نقطه توازن میان کیفیت فارسی، سرعت پاسخ و مصرف منابع
قرار گیرد.
بر همین اساس، چند مدل در اندازه‌های متفاوت به‌صورت تجربی
روی سخت‌افزار یکسان مقایسه شده‌اند و نتایج آن‌ها در فصل
آزمایش‌ها ارائه خواهد شد.
% =========================================================
% 3-6
% =========================================================
\section{پژوهش‌های مرتبط با تبدیل متن به گفتار فارسی}

در مقایسه با زبان‌هایی مانند انگلیسی، توسعه سامانه‌های
تبدیل متن به گفتار برای زبان فارسی مدت زیادی با محدودیت
مجموعه‌داده‌های بزرگ و باکیفیت مواجه بوده است.

مدل‌های عصبی جدید \lr{TTS} معمولاً برای یادگیری تلفظ طبیعی،
ریتم، آهنگ و ویژگی‌های گوینده به داده صوتی دقیق و هم‌تراز
با متن نیاز دارند.
در نتیجه، کیفیت و حجم داده آموزشی نقش بسیار مهمی در عملکرد
این مدل‌ها دارد.

در سال‌های اخیر، چند مجموعه داده فارسی با هدف توسعه
سامانه‌های \lr{TTS} منتشر شده‌اند و پژوهشگران از معماری‌هایی مانند
Tacotron 2، HiFi-GAN و \lr{VITS} برای تولید گفتار فارسی استفاده
کرده‌اند.

در ادامه، مهم‌ترین پژوهش‌ها و منابع مرتبط با موضوع پروژه حاضر
بررسی می‌شوند.


% ---------------------------------------------------------
\subsection{مجموعه داده ArmanTTS}

یکی از مجموعه‌های فارسی طراحی‌شده برای تبدیل متن به گفتار،
\lr{ArmanTTS}
است که توسط
\lr{Shamgholi et al.}
معرفی شده است
\cite{shamgholi2023armantts}.

این مجموعه یک پایگاه داده تک‌گوینده
\LTRfootnote{Single-Speaker}
برای زبان فارسی است.

پژوهشگران علاوه بر معرفی مجموعه داده، یک سامانه تبدیل متن
به گفتار مبتنی بر
\lr{Tacotron 2}
و
\lr{HiFi-GAN}
پیاده‌سازی کردند.

در این سامانه، ورودی متنی ابتدا به نمایش واجی تبدیل می‌شود
و سپس Tacotron 2 برای تولید نمایش میانی گفتار و HiFi-GAN
برای تولید شکل موج صوتی مورد استفاده قرار می‌گیرد.

ارزیابی انسانی انجام‌شده در این پژوهش نشان داد که امتیاز
\lr{MOS}
برای گفتار طبیعی حدود 4.0، برای خروجی واکودر حدود 3.87
و برای گفتار تولیدشده توسط کل سامانه حدود 2.98 بوده است.

این اختلاف نشان می‌دهد که کیفیت واکودر به‌تنهایی بالا بوده،
اما بخش تبدیل متن به بازنمایی صوتی هنوز عامل محدودکننده
کیفیت نهایی بوده است.

\textbf{ارتباط با پروژه حاضر:}
این پژوهش اهمیت تفکیک منابع مختلف افت کیفیت را نشان می‌دهد.
در یک سامانه \lr{TTS}، کیفیت نهایی تنها به مرحله تولید waveform
وابسته نیست و بخش‌هایی مانند نرمال‌سازی متن، مدل آکوستیکی،
هم‌ترازی و تلفظ نیز می‌توانند بر نتیجه اثر بگذارند.


% ---------------------------------------------------------
\subsection{مجموعه \lr{DeepMine-Multi-TTS}}

یکی دیگر از تلاش‌ها برای توسعه منابع فارسی،
مجموعه
\lr{DeepMine-Multi-TTS}
است
\cite{adibian2023deepminemultitts}.

برخلاف ArmanTTS که یک مجموعه تک‌گوینده است،
\lr{DeepMine-Multi-TTS} با هدف توسعه مدل‌های
\textbf{چندگوینده}
\LTRfootnote{Multi-Speaker}
طراحی شده است.

این مجموعه شامل حدود 120 ساعت گفتار باکیفیت از 67 گوینده
فارسی‌زبان است.

پژوهشگران با استفاده از این مجموعه چند سامانه تولید گفتار
را آموزش داده و کیفیت خروجی را با معیار \lr{MOS} ارزیابی کردند.

طبق نتایج گزارش‌شده، دو مدل چندگوینده مورد آزمایش به
\lr{MOS}های تقریباً 3.94 و 4.12 دست یافتند.

این نتایج نشان می‌دهند که وجود تنوع گویندگان در مجموعه داده
می‌تواند امکان آموزش مدل‌های انعطاف‌پذیرتر را فراهم کند.

\textbf{ارتباط با پروژه حاضر:}
اگرچه هدف پروژه حاضر شبیه‌سازی چند گوینده نیست، این پژوهش
نشان می‌دهد که کیفیت و تنوع داده فارسی می‌تواند تأثیر مستقیمی
بر طبیعی‌بودن خروجی داشته باشد.


% ---------------------------------------------------------
\subsection{مجموعه ManaTTS}

یکی از منابع جدید و مهم برای پژوهش \lr{TTS} فارسی،
\lr{ManaTTS}
است که توسط
\lr{Qharabagh et al.}
معرفی شده است
\cite{qharabagh2025manatts}.

هدف این پژوهش تنها انتشار یک مجموعه داده نبود، بلکه یک
فرآیند کامل برای ساخت مجموعه‌های \lr{TTS} در زبان‌های کم‌منبع
ارائه شد.

نسخه معرفی‌شده در مقاله شامل حدود 86 ساعت صوت فارسی
تک‌گوینده با نرخ نمونه‌برداری 44.1 کیلوهرتز بود.

خط لوله ایجاد داده شامل مؤلفه‌هایی مانند:

\begin{itemize}

    \item قطعه‌بندی متن؛
    \item قطعه‌بندی محدودشده صوت؛
    \item هم‌ترازی صوت و متن؛
    \item کنترل کیفیت نمونه‌ها؛
    \item و آماده‌سازی داده برای آموزش \lr{TTS}
\end{itemize}

بود.

برای بررسی کیفیت مجموعه داده، پژوهشگران یک مدل مبتنی بر
Tacotron 2 آموزش دادند.

گفتار تولیدشده توسط این مدل به
\lr{MOS}
حدود 3.76 دست یافت، در حالی که \lr{MOS} مربوط به waveform طبیعی
حدود 4.01 گزارش شد.

فاصله نسبتاً محدود میان گفتار سنتزشده و گفتار طبیعی نشان
می‌دهد که مجموعه ایجادشده از کیفیت مناسبی برای آموزش \lr{TTS}
برخوردار بوده است.

\textbf{ارتباط با پروژه حاضر:}
این مطالعه اهمیت داده آموزشی و فرآیند پاک‌سازی آن را نشان
می‌دهد.
حتی با ثابت‌بودن معماری، تغییر کیفیت داده می‌تواند تفاوت
زیادی در خروجی ایجاد کند.


% ---------------------------------------------------------
\subsection{گسترش منابع فارسی با ParsVoice}

یکی از جدیدترین منابع در حوزه \lr{TTS} فارسی،
\lr{ParsVoice}
است
\cite{ranjbar2025parsvoice}.

این مجموعه با استفاده از کتاب‌های صوتی فارسی ایجاد شده و
هدف آن فراهم‌کردن یک پایگاه داده بسیار بزرگ چندگوینده برای
تولید گفتار است.

پژوهشگران یک خط لوله خودکار برای تبدیل کتاب صوتی خام به
نمونه‌های مناسب \lr{TTS} توسعه دادند.
این خط لوله شامل تشخیص کامل‌بودن جمله، بهینه‌سازی مرزهای
صوت و متن و ارزیابی چندبعدی کیفیت نمونه‌ها بود.

در مجموع حدود 3526 ساعت صوت پردازش شد و پس از اعمال
فیلترهای کیفیت، حدود 1804 ساعت گفتار باکیفیت از بیش از
470 گوینده به‌عنوان مجموعه اصلی ارائه شد.

این حجم داده در مقایسه با مجموعه‌های قدیمی‌تر فارسی افزایش
قابل توجهی محسوب می‌شود.

\textbf{ارتباط با پروژه حاضر:}
ظهور مجموعه‌هایی مانند ParsVoice نشان می‌دهد که یکی از
محدودیت‌های تاریخی \lr{TTS} فارسی، یعنی کمبود داده باکیفیت،
در حال کاهش است.
این موضوع می‌تواند در آینده امکان آموزش مدل‌های فارسی
قدرتمندتر و متنوع‌تر را فراهم کند.


% ---------------------------------------------------------
\subsection{مدل‌های چندزبانه \lr{MMS} برای فارسی}

پروژه
\lr{Massively Multilingual \lr{Speech} (\lr{MMS})}
یکی از تلاش‌های گسترده برای توسعه فناوری گفتار در تعداد
زیادی از زبان‌ها است
\cite{pratap2023mms}.

این پروژه علاوه بر تشخیص گفتار، شامل مدل‌های
تبدیل متن به گفتار برای تعداد زیادی از زبان‌ها نیز می‌شود.

برای زبان فارسی، یک checkpoint اختصاصی با شناسه:

\begin{center}
\lr{facebook/mms-tts-fas}
\end{center}

ارائه شده است.

این مدل از معماری
\lr{VITS}
استفاده می‌کند و متن فارسی را مستقیماً به شکل موج گفتاری
تبدیل می‌کند.

مزیت اصلی این رویکرد آن است که بدون نیاز به آموزش یک مدل
جدید از ابتدا، امکان استفاده از یک مدل آماده برای زبان فارسی
فراهم می‌شود.

\textbf{ارتباط با پروژه حاضر:}
مدل \lr{MMS} یکی از سه خانواده اصلی \lr{TTS} مورد ارزیابی در پروژه
حاضر است.
از آنجا که این مدل بخشی از یک پروژه بزرگ چندزبانه است،
ارزیابی تجربی آن در کنار مدل‌های فارسی دیگر می‌تواند نشان
دهد که آیا یک مدل چندزبانه عمومی می‌تواند در کاربرد محلی
فارسی عملکرد مناسبی داشته باشد.


% ---------------------------------------------------------
\subsection{مدل‌های مبتنی بر \lr{VITS}}

بخش قابل توجهی از مدل‌های جدید \lr{TTS} فارسی بر پایه معماری
\lr{VITS} یا مشتقات آن ساخته شده‌اند.

همان‌طور که در فصل مفاهیم نظری توضیح داده شد، \lr{VITS} یک
معماری انتهابه‌انتها برای تبدیل متن به waveform است که
مدل‌سازی نهفته، جریان‌های نرمال‌ساز و آموزش خصمانه را
ترکیب می‌کند
\cite{kim2021vits}.

این معماری به دلیل کیفیت صوت مناسب و حذف نیاز به یک
واکوُدر مستقل در زمان استنتاج، برای سامانه‌های عملی مناسب
است.

در پروژه حاضر، دو مدل از مدل‌های مورد آزمایش مستقیماً در
خانواده \lr{VITS} قرار می‌گیرند:

\begin{itemize}

    \item \lr{Kamtera\_VITS}
    \item \lr{VITS\_MMS\_FAS}

\end{itemize}

استفاده از دو مدل مبتنی بر یک خانواده معماری ولی با
داده و فرآیند آموزش متفاوت، امکان بررسی اثر داده و تنظیم
مدل بر کیفیت نهایی را فراهم می‌کند.


% ---------------------------------------------------------
\subsection{سامانه \lr{Piper} و تولید گفتار محلی}

در کنار مدل‌های عمومی \lr{VITS}، ابزار
\lr{Piper}
با تمرکز ویژه بر تولید گفتار سریع و محلی توسعه یافته است.

\lr{Piper} از مدل‌های مبتنی بر \lr{VITS} استفاده کرده و آن‌ها را
برای اجرای سریع با
\lr{ONNX Runtime}
آماده می‌کند.

در مجموعه رسمی صداهای \lr{Piper}، مدل‌های فارسی با شناسه
\lr{fa\_IR}
ارائه شده‌اند.

از جمله صداهای فارسی موجود می‌توان به مدل‌هایی مانند:

\begin{center}
\lr{amir}
\end{center}

و:

\begin{center}
\lr{gyro}
\end{center}

اشاره کرد.

تمرکز \lr{Piper} بر اجرای سبک و محلی باعث شده است این ابزار برای
دستیارهای صوتی آفلاین، دستگاه‌های لبه و سامانه‌های تعاملی
مناسب باشد.

\textbf{ارتباط با پروژه حاضر:}
مدل \lr{Piper} فارسی یکی از گزینه‌های اصلی بنچمارک \lr{TTS} این پروژه
است.
هدف از بررسی آن این است که مشخص شود آیا مدل بهینه‌شده برای
اجرای محلی می‌تواند در کنار سرعت مناسب، کیفیت صوت قابل
قبولی نیز ارائه کند.


% ---------------------------------------------------------
\subsection{کیفیت ادراکی و معیار \lr{MOS}}

بخش بزرگی از پژوهش‌های \lr{TTS} هنوز از
\lr{Mean Opinion Score}
به‌عنوان معیار اصلی طبیعی‌بودن گفتار استفاده می‌کنند.

مطالعات ArmanTTS، \lr{DeepMine-Multi-TTS} و ManaTTS همگی از
ارزیابی انسانی \lr{MOS} برای مقایسه کیفیت خروجی استفاده کرده‌اند
\cite{shamgholi2023armantts,
      adibian2023deepminemultitts,
      qharabagh2025manatts}.

این امر نشان می‌دهد که با وجود پیشرفت معیارهای خودکار،
ادراک انسانی همچنان یکی از منابع اصلی ارزیابی کیفیت
گفتار مصنوعی است.

با این حال، \lr{MOS} دارای محدودیت‌هایی نیز هست.
نتایج می‌توانند به تعداد شنوندگان، کیفیت هدفون، محیط شنیدن،
دستورالعمل آزمایش و تفاوت سلیقه افراد وابسته باشند.

به همین دلیل، استفاده از معیارهای خودکار در کنار \lr{MOS} می‌تواند
برای آزمایش‌های بزرگ‌تر مفید باشد.


% ---------------------------------------------------------
\subsection{ارزیابی خودکار کیفیت گفتار}

ارزیابی دستی هزاران نمونه صوتی هزینه زمانی زیادی دارد.
به همین دلیل، مدل‌هایی مانند
\lr{NISQA}
برای تخمین خودکار کیفیت گفتار توسعه یافته‌اند
\cite{mittag2021nisqa}.

\lr{NISQA} علاوه بر امتیاز کلی کیفیت، ابعادی مانند نویز،
ناپیوستگی، coloration و loudness را نیز برآورد می‌کند.

در پروژه حاضر این ویژگی اهمیت زیادی دارد، زیرا تعداد نمونه‌های
ارزیابی \lr{TTS} نسبتاً زیاد است و ارزیابی تمام آن‌ها توسط
شنوندگان انسانی عملی نیست.

استفاده از \lr{NISQA} امکان مقایسه استانداردتر تعداد زیادی فایل
خروجی را فراهم می‌کند.

با این حال، مقدار تولیدشده توسط \lr{NISQA} باید به‌عنوان
\textbf{تخمین مدل}
در نظر گرفته شود و نباید کاملاً معادل \lr{MOS} انسانی تلقی شود.


% ---------------------------------------------------------
\subsection{ارزیابی قابل فهم بودن خروجی}

کیفیت طبیعی صوت و قابل‌فهم‌بودن آن دو مفهوم مرتبط اما
متفاوت هستند.

ممکن است یک فایل صوتی از نظر نویز و پیوستگی کیفیت مناسبی داشته
باشد، اما برخی واژه‌ها به‌درستی تلفظ نشوند.

به همین دلیل، برخی ارزیابی‌ها از \lr{ASR} برای تبدیل مجدد صوت
تولیدشده به متن استفاده می‌کنند.

در این روش:

\begin{center}
Reference \lr{Text}
$\longrightarrow$
\lr{TTS}
$\longrightarrow$
Synthetic \lr{Speech}
$\longrightarrow$
\lr{ASR}
$\longrightarrow$
Recognized \lr{Text}
\end{center}

و سپس متن مرجع و متن بازشناسی‌شده با معیارهای
\lr{WER}
و
\lr{CER}
مقایسه می‌شوند.

این روش امکان ارزیابی خودکار intelligibility را فراهم می‌کند،
اما نتیجه به کیفیت \lr{ASR} ارزیاب نیز وابسته است.

بنابراین استفاده از این معیار در کنار ارزیابی ادراکی
مناسب‌تر از استفاده مستقل از آن است.


% ---------------------------------------------------------
\subsection{سرعت تولید و کاربرد مکالمه‌ای}

یکی از تفاوت‌های مهم \lr{TTS} در یک چت‌بات با کاربردهایی مانند
تولید کتاب صوتی، اهمیت بسیار بیشتر تأخیر است.

در کتاب صوتی ممکن است تولید یک جمله چند ثانیه طول بکشد و
همچنان قابل قبول باشد.
اما در گفت‌وگو، کاربر انتظار دارد پاسخ تقریباً بلافاصله
پس از تولید متن آغاز شود.

به همین دلیل، پژوهش درباره مدل‌های محلی مانند \lr{Piper} اهمیت
ویژه‌ای برای کاربردهای مکالمه‌ای دارد.

برای چنین کاربردی باید حداقل سه بعد هم‌زمان بررسی شوند:

\begin{equation}
Q_{\\mathrm{TTS}}
=
f(
Q_{\mathrm{perceptual}},
I_{\mathrm{intelligibility}},
T_{\mathrm{synthesis}}
)
\end{equation}

که در آن:

\begin{itemize}

    \item $Q_{\mathrm{perceptual}}$ کیفیت ادراکی گفتار؛
    \item $I_{\mathrm{intelligibility}}$ میزان قابل فهم بودن؛
    \item و $T_{\mathrm{synthesis}}$ هزینه زمانی تولید است.

\end{itemize}

ازاین‌رو، انتخاب مدل \lr{TTS} برای چت‌بات باید بر اساس توازن
میان این سه عامل انجام شود.


% ---------------------------------------------------------
\subsection{جمع‌بندی پژوهش‌های \lr{TTS} فارسی}

مرور پژوهش‌های این حوزه نشان می‌دهد که یکی از مهم‌ترین
چالش‌های تاریخی تبدیل متن به گفتار فارسی کمبود داده باکیفیت
بوده است.

مجموعه‌هایی مانند ArmanTTS منابع تک‌گوینده اولیه‌تری را
فراهم کردند، در حالی که \lr{DeepMine-Multi-TTS} تنوع چندگوینده
بیشتری ارائه داد.
پژوهش ManaTTS علاوه بر مجموعه داده، یک خط لوله کامل برای
ساخت داده \lr{TTS} در زبان‌های کم‌منبع معرفی کرد.
در سال‌های جدیدتر نیز ParsVoice حجم داده فارسی را به مقیاس
بسیار بزرگ‌تری افزایش داده است
\cite{shamgholi2023armantts,
      adibian2023deepminemultitts,
      qharabagh2025manatts,
      ranjbar2025parsvoice}.

از نظر معماری، مدل‌های مبتنی بر \lr{VITS} و پروژه چندزبانه \lr{MMS}
امکان تولید گفتار فارسی با مدل‌های عصبی انتهابه‌انتها را
فراهم کرده‌اند.

از سوی دیگر، \lr{Piper} نمونه‌ای از تلاش برای انتقال این مدل‌ها
به محیط‌های محلی و کم‌تأخیر است.

این مطالعات نشان می‌دهند که برای پروژه حاضر، انتخاب مدل \lr{TTS}
نباید تنها بر اساس طبیعی‌بودن صوت انجام شود.
قابل فهم بودن گفتار، زمان تولید و امکان اجرای محلی نیز باید
در تصمیم نهایی مورد توجه قرار گیرند.

به همین دلیل، مدل‌های \lr{TTS} پروژه حاضر با ترکیبی از معیارهای
کیفیت ادراکی، intelligibility و عملکرد اجرایی مقایسه
شده‌اند.
% =========================================================
% 3-7
% =========================================================
\section{سامانه‌های گفت‌وگوی صوتی یکپارچه}

سامانه‌های گفت‌وگوی صوتی را می‌توان به دو گروه کلی تقسیم کرد:
سامانه‌های زنجیره‌ای و مدل‌های انتهابه‌انتها.

در معماری زنجیره‌ای، گفتار کاربر ابتدا توسط \lr{ASR} به متن تبدیل
می‌شود، سپس مدل زبانی پاسخ را تولید می‌کند و در نهایت \lr{TTS}
پاسخ متنی را به گفتار تبدیل می‌کند:

\begin{center}
\lr{Speech}
$\longrightarrow$
\lr{ASR}
$\longrightarrow$
\lr{LLM}
$\longrightarrow$
\lr{TTS}
$\longrightarrow$
\lr{Speech}
\end{center}

این معماری به دلیل ماژولار بودن، امکان انتخاب و جایگزینی مستقل
هر مؤلفه را فراهم می‌کند.
در مقابل، یکی از محدودیت‌های آن انتقال خطا میان مراحل و افزایش
تأخیر ناشی از اجرای چند مدل پشت سر هم است
\cite{cui2024speechlmsurvey}.


% ---------------------------------------------------------
\subsection{حرکت به سمت مدل‌های \lr{Speech} Language \lr{Model}}

در سال‌های اخیر، پژوهش‌ها به سمت مدل‌هایی حرکت کرده‌اند که
بتوانند صوت را مستقیماً دریافت کرده و بدون تبدیل اجباری به متن،
آن را پردازش کنند.

یکی از نمونه‌های اولیه،
\lr{SpeechGPT}
است
\cite{zhang2023speechgpt}.

در این مدل، بازنمایی‌های گسسته گفتار در کنار مدل زبانی استفاده
می‌شوند تا مدل بتواند ورودی و خروجی صوتی را در یک چارچوب واحد
پردازش کند.

هدف این رویکرد کاهش فاصله میان مدالیته گفتار و متن و ایجاد
قابلیت مکالمه مستقیم‌تر است.

با این حال، چنین مدل‌هایی معمولاً نیازمند منابع محاسباتی و
داده آموزشی بسیار بیشتری نسبت به یک سامانه ماژولار هستند.


% ---------------------------------------------------------
\subsection{AudioPaLM}

مدل
\lr{AudioPaLM}
نیز با هدف ترکیب توانایی مدل‌های زبانی متنی با مدل‌های گفتاری
توسعه یافته است
\cite{rubenstein2023audiopalm}.

این مدل می‌تواند هم ورودی متنی و هم صوتی را پردازش کند و
وظایفی مانند تشخیص گفتار و ترجمه گفتار را در یک معماری مشترک
انجام دهد.

یکی از ویژگی‌های مهم این مدل، حفظ بخشی از اطلاعات فرازبانی
مانند ویژگی‌های گوینده و آهنگ گفتار است؛ اطلاعاتی که در یک
خط لوله صرفاً متنی ممکن است از بین بروند.


% ---------------------------------------------------------
\subsection{SALMONN}

مدل
\lr{SALMONN}
نمونه دیگری از مدل‌های چندوجهی صوت و زبان است
\cite{tang2023salmonn}.

در این معماری، رمزگذارهای صوت و گفتار به یک مدل زبانی متنی
متصل می‌شوند تا مدل بتواند علاوه بر گفتار، رویدادهای صوتی و
موسیقی را نیز پردازش کند.

این مدل نشان می‌دهد که مدل‌های زبانی جدید می‌توانند به‌تدریج
از پردازش صرف متن فراتر رفته و نقش هسته مرکزی یک سامانه شنیداری
عمومی را ایفا کنند.


% ---------------------------------------------------------
\subsection{مقایسه رویکرد زنجیره‌ای و انتهابه‌انتها}

مدل‌های گفتاری انتهابه‌انتها مزایایی مانند کاهش تبدیل میان
گفتار و متن و امکان حفظ اطلاعات صوتی بیشتری دارند.

در مقابل، معماری زنجیره‌ای چند مزیت عملی مهم دارد:

\begin{itemize}

    \item امکان انتخاب مستقل بهترین مدل \lr{ASR}، \lr{LLM} و \lr{TTS}؛
    \item امکان ارزیابی و اشکال‌زدایی جداگانه هر بخش؛
    \item امکان تعویض یک مؤلفه بدون آموزش مجدد کل سامانه؛
    \item و امکان استفاده از مدل‌های سبک‌تر برای اجرای محلی.
\end{itemize}

به همین دلیل، با وجود پیشرفت مدل‌های \lr{Speech} Language \lr{Model}،
معماری زنجیره‌ای همچنان برای سامانه‌هایی با محدودیت منابع
یک انتخاب عملی محسوب می‌شود.


% ---------------------------------------------------------
\subsection{جایگاه پروژه حاضر}

پروژه حاضر از معماری ماژولار:

\begin{center}
\lr{ASR}
$\longrightarrow$
\lr{LLM}
$\longrightarrow$
\lr{TTS}
\end{center}

استفاده می‌کند.

این انتخاب با هدف اجرای کامل سامانه روی سخت‌افزار شخصی،
امکان مقایسه چند مدل مستقل و کنترل بهتر مصرف منابع انجام شده
است.

در مقایسه با مدل‌های بزرگ گفتار-زبان، این معماری امکان
بهینه‌سازی مستقل هر مرحله را فراهم می‌کند، هرچند چالش‌هایی
مانند انتشار خطا و مجموع تأخیر مراحل همچنان باقی می‌مانند.

به همین دلیل، در پروژه حاضر علاوه بر ارزیابی مستقل اجزا،
ارزیابی انتهابه‌انتهای سامانه نیز اهمیت دارد.
% =========================================================
% 3-8
% =========================================================
\section{مقایسه پژوهش‌های پیشین}

مرور پژوهش‌های انجام‌شده نشان می‌دهد که در هر یک از سه مؤلفه
اصلی سامانه گفت‌وگوی صوتی، رویکردهای متفاوتی برای ایجاد توازن
میان کیفیت و هزینه محاسباتی ارائه شده است.

در حوزه تشخیص گفتار، مدل‌های چندزبانه مانند
\lr{XLSR}
و
\lr{Whisper}
امکان انتقال دانش میان زبان‌ها را فراهم کرده‌اند، در حالی که
مدل‌های اختصاصی‌تر مانند
\lr{FastConformer}
می‌توانند در برخی شرایط برای زبان فارسی عملکرد رقابتی‌تری
داشته باشند.
پژوهش‌های مرتبط همچنین نشان داده‌اند که عواملی مانند کیفیت
داده، نوع گفتار، لهجه، نویز و نحوه قطعه‌بندی صوت می‌توانند
تأثیر قابل توجهی بر عملکرد \lr{ASR} داشته باشند
\cite{conneau2020xlsr,radford2022whisper,sedghiyeh2025psrb}.

در حوزه مدل‌های زبانی، مطالعات فارسی نشان می‌دهند که مدل
بزرگ‌تر الزاماً بهترین گزینه برای تمام وظایف نیست.
مدل‌های عمومی چندزبانه از دانش گسترده‌تری برخوردارند، اما
مدل‌های تنظیم‌شده فارسی ممکن است در برخی کاربردهای بومی
عملکرد بهتری ارائه کنند
\cite{abaskohi2024persianllm,hosseinbeigi2025persianllm}.

همچنین، برای اجرای محلی، عواملی مانند زمان پاسخ، سرعت تولید
توکن و مصرف حافظه اهمیت زیادی دارند.
روش‌هایی مانند کوانتیزه‌سازی نیز می‌توانند امکان اجرای مدل‌های
بزرگ‌تر روی سخت‌افزار محدود را فراهم کنند
\cite{frantar2022gptq,lin2023awq}.

در حوزه تبدیل متن به گفتار، پژوهش‌های فارسی نشان داده‌اند که
کیفیت مجموعه‌داده نقش اساسی در طبیعی‌بودن خروجی دارد.
در کنار کیفیت ادراکی، سرعت تولید و قابل‌فهم‌بودن گفتار نیز
برای کاربردهای تعاملی اهمیت دارند
\cite{shamgholi2023armantts,qharabagh2025manatts}.

در سطح سامانه کامل نیز دو رویکرد اصلی مشاهده می‌شود.
مدل‌های جدید گفتار-زبان تلاش می‌کنند پردازش صوت و زبان را
در یک معماری مشترک انجام دهند، در حالی که معماری‌های زنجیره‌ای
هنوز به دلیل ماژولار بودن و امکان کنترل مستقل مؤلفه‌ها در
کاربردهای عملی مورد استفاده قرار می‌گیرند
\cite{zhang2023speechgpt,rubenstein2023audiopalm}.

به‌طور خلاصه، می‌توان تفاوت رویکردهای موجود را چنین بیان کرد:

\begin{itemize}

    \item مدل‌های بزرگ معمولاً توانایی عمومی بیشتری دارند،
    اما هزینه محاسباتی بالاتری ایجاد می‌کنند؛

    \item مدل‌های اختصاصی زبان هدف می‌توانند با اندازه کمتر
    عملکرد رقابتی ارائه دهند؛

    \item ارزیابی مستقل یک مؤلفه برای انتخاب آن در یک سامانه
    کامل کافی نیست؛

    \item و در کاربردهای محلی، معیارهایی مانند تأخیر و مصرف
    حافظه باید در کنار کیفیت در نظر گرفته شوند.

\end{itemize}

بنابراین، ادبیات موجود از یک پاسخ واحد برای انتخاب بهترین
مدل \lr{ASR}، \lr{LLM} یا \lr{TTS} حمایت نمی‌کند.
انتخاب مناسب به زبان، داده، سخت‌افزار و کاربرد نهایی وابسته
است.
این مسئله زمینه اصلی برای تعریف شکاف پژوهشی پروژه حاضر را
فراهم می‌کند.

% =========================================================
% 3-9
% =========================================================
\section{شکاف پژوهشی و جایگاه پروژه حاضر}

مرور پژوهش‌های پیشین نشان می‌دهد که بخش زیادی از مطالعات
مرتبط، هر یک تنها یکی از مؤلفه‌های سامانه گفت‌وگوی صوتی را
به‌صورت مستقل بررسی کرده‌اند.

در حوزه تشخیص گفتار، تمرکز اصلی بر کاهش
\lr{WER}
و بهبود مقاومت مدل در برابر تفاوت گوینده، لهجه و شرایط صوتی
بوده است.
در حوزه مدل‌های زبانی نیز پژوهش‌ها بیشتر بر کیفیت پاسخ،
درک زبان فارسی و ارزیابی دانش زبانی و فرهنگی تمرکز داشته‌اند.
در مطالعات \lr{TTS} نیز طبیعی‌بودن گفتار، کیفیت ادراکی و
قابل‌فهم‌بودن خروجی از مهم‌ترین معیارهای ارزیابی بوده‌اند.

با این حال، در یک چت‌بات صوتی واقعی این مؤلفه‌ها به‌صورت
مستقل عمل نمی‌کنند، بلکه در قالب یک خط لوله متوالی به یکدیگر
وابسته هستند:

\begin{center}
\lr{ASR}
$\longrightarrow$
\lr{LLM}
$\longrightarrow$
\lr{TTS}
\end{center}

بنابراین، انتخاب بهترین مدل مستقل در هر بخش الزاماً به بهترین
سامانه نهایی منجر نمی‌شود.
برای مثال، یک مدل \lr{ASR} ممکن است دقت بسیار خوبی داشته باشد،
اما مصرف حافظه یا زمان استنتاج آن به اندازه‌ای زیاد باشد که
اجرای هم‌زمان آن با مدل زبانی و \lr{TTS} روی یک رایانه شخصی
دشوار شود.

به همین شکل، یک مدل زبانی بزرگ‌تر ممکن است پاسخ‌های باکیفیت‌تری
تولید کند، اما تأخیر بیشتر آن تجربه مکالمه را کاهش دهد.
در بخش \lr{TTS} نیز طبیعی‌بودن صوت باید در کنار سرعت تولید بررسی شود.

یکی دیگر از شکاف‌های مهم، کمبود ارزیابی‌هایی است که
\textbf{کل خط لوله فارسی را تحت محدودیت اجرای محلی}
بررسی کنند.

پروژه حاضر تلاش می‌کند این مسئله را با ارزیابی چندمرحله‌ای
سامانه برطرف کند.
در این پروژه، چند مدل مختلف برای هر یک از سه مؤلفه اصلی
مقایسه شده‌اند و معیارهای ارزیابی تنها به کیفیت محدود نیستند.

برای \lr{ASR} معیارهایی مانند:

\begin{itemize}
    \item \lr{WER} و \lr{CER}؛
    \item \lr{RTF}؛
    \item زمان تأخیر؛
    \item و مصرف حافظه
\end{itemize}

در نظر گرفته شده‌اند.

در بخش مدل زبانی نیز معیارهایی مانند کیفیت معنایی پاسخ،
زمان پاسخ، سرعت تولید توکن و مصرف حافظه بررسی می‌شوند.

در بخش \lr{TTS} نیز کیفیت ادراکی، قابل‌فهم‌بودن گفتار،
زمان تولید و ضریب زمان واقعی مورد مقایسه قرار می‌گیرند.

در نتیجه، هدف اصلی پروژه حاضر پیدا کردن
\textbf{بهترین توازن عملی میان کیفیت و کارایی}
برای ساخت یک چت‌بات صوتی فارسی محلی است، نه صرفاً انتخاب
مدلی که در یک معیار منفرد بهترین مقدار را داشته باشد.

ویژگی دیگر پروژه، معماری ماژولار آن است.
این معماری امکان جایگزینی مستقل مدل‌های \lr{ASR}، \lr{LLM} و \lr{TTS} را
فراهم می‌کند و اجازه می‌دهد هر مؤلفه به‌صورت جداگانه و سپس
در قالب سامانه کامل ارزیابی شود.

همچنین، مؤلفه‌هایی مانند
\lr{VAD}،
پردازش نویز،
مدیریت بافت مکالمه
و ارزیابی تأخیر انتهابه‌انتها
به‌عنوان بخش‌های تکمیلی سامانه در نظر گرفته شده‌اند.

در مجموع، جایگاه پروژه حاضر را می‌توان در سه محور خلاصه کرد:

\begin{itemize}

    \item مقایسه تجربی چند مدل فارسی یا چندزبانه در هر یک از
    بخش‌های \lr{ASR}، \lr{LLM} و \lr{TTS}؛

    \item انتخاب مدل‌ها بر اساس توازن میان کیفیت، تأخیر و
    مصرف منابع؛

    \item و یکپارچه‌سازی مدل‌های منتخب در یک سامانه گفت‌وگوی
    صوتی فارسی با قابلیت اجرای محلی.

\end{itemize}

بر این اساس، پروژه حاضر در نقطه اتصال پژوهش‌های مستقل
پردازش گفتار، مدل‌های زبانی و تولید گفتار قرار می‌گیرد و
تمرکز اصلی آن تبدیل این اجزای مستقل به یک سامانه عملی و
قابل ارزیابی است.
% =========================================================
% 3-10
% =========================================================
\section{جمع‌بندی}

در این فصل، پژوهش‌های مرتبط با تشخیص گفتار فارسی،
مدل‌های چندزبانه و کم‌منبع، مدل‌های زبانی فارسی،
اجرای محلی مدل‌های زبانی، تبدیل متن به گفتار فارسی
و سامانه‌های گفت‌وگوی صوتی بررسی شدند.

مطالعات پیشین نشان می‌دهند که کیفیت داده، سازگاری مدل با
زبان فارسی و نوع کاربرد نقش مهمی در عملکرد نهایی دارند.
همچنین، مدل بزرگ‌تر الزاماً بهترین انتخاب برای یک سامانه
محلی نیست و معیارهایی مانند سرعت و مصرف منابع باید در کنار
کیفیت در نظر گرفته شوند.

بررسی پژوهش‌های مرتبط همچنین نشان داد که بخش زیادی از
مطالعات، مؤلفه‌های \lr{ASR}، \lr{LLM} و \lr{TTS} را به‌صورت مستقل بررسی
کرده‌اند.
در مقابل، پروژه حاضر بر مقایسه، انتخاب و یکپارچه‌سازی این
مؤلفه‌ها در قالب یک چت‌بات صوتی فارسی محلی تمرکز دارد.

در فصل بعد، معماری کلی سامانه پیشنهادی، نحوه ارتباط ماژول‌ها
و تصمیم‌های طراحی پروژه تشریح خواهد شد.
